% The Kotlin chapter of the WFDB-x user's manual. It follows the sections of the Python
% chapter (python.tex), in the same order. Each \apientry of the section "The high-level
% API" documents one public high-level function of kotlin/src/commonMain/kotlin/wfdbx
% (Record.kt, Check.kt, Metadata.kt): its signature listing holds the declarations of the
% extension function of Wfdb (all its overloads), copied from the code without their
% bodies, and its params name every parameter of these declarations, in code order. Each
% function is also a top-level function with the same parameters (the default instance).
% The listings of style "example" were compiled and run against the library.
\chapter{Kotlin}
\label{chap:kotlin}
\chapterlanguage{kotlin}

The Kotlin implementation of WFDB-x is the Kotlin Multiplatform library \code{wfdbx}, with
the Gradle coordinates \code{io.github.alphanumericslab:wfdbx:1.2.0}, in the directory
\file{kotlin} of the repository. It is written in common Kotlin code, which runs on every
JVM, on Android and on iOS, with the Kotlin standard library as its only dependency, and
never calls the C library. It is meant for mobile applications and JVM programs that read
and write WFDB records on the device, with the behavior of the other implementations of
WFDB-x, without native libraries, JNI or another runtime.

It offers WFDB-x in three ways:
\begin{itemize}
\item \textbf{The high-level API} (\cref{sec:kotlin-api}): functions that read and write
  whole records, headers, annotation files and header metadata in one call, with named
  arguments and default values; the samples are an \code{IntArray} (and a
  \code{DoubleArray} in physical units), frame by frame, and errors throw a
  \code{WfdbException}. They read signal files in large blocks, and are the fastest way
  to read a record. They are an addition of WFDB-x: the C package has no equivalent. They
  are the recommended interface for new code and apps.
\item \textbf{The WFDB library functions} under their C names (\code{isigopen},
  \code{getvec}, \code{getann}, \code{putann}, \ldots), with the behavior of the C library
  (\cref{sec:kotlin-library}). They are documented once, for all the languages, in
  \cref{chap:library}.
\item \textbf{The command-line tools} (\code{rdsamp}, \code{rdann}, \code{wrsamp},
  \code{xform} and 29 others), with the options, output and exit status of the C tools,
  run on the JVM through the launcher \file{kotlin/bin/wfdbx} (\cref{sec:kotlin-tools}).
  They are documented in \cref{chap:tools}.
\end{itemize}
The high-level functions return exactly the samples and annotations that the library
functions (and the C library) give.

% ==========================================================================================
\section{Installation and requirements}
\label{sec:kotlin-install}

\subsection{Requirements}

The library is built with Gradle 8.14.3 (the Gradle wrapper of the repository,
\file{kotlin/gradlew}, which downloads itself) and Kotlin 2.1.21, and needs a JDK 11 or
later to build. It has no dependency besides the Kotlin standard library (and
\code{kotlin-test} for the tests); the FLAC formats 508, 516 and 524 are read and written
by WFDB-x's own codec. The targets are:
\begin{itemize}
\item \code{jvm()}: bytecode for Java 11, which runs on every current JVM and on Android
  (API level 26 or later; the library uses \code{java.nio.file});
\item \code{iosArm64()}, \code{iosSimulatorArm64()} and \code{iosX64()}: declared, with a
  small platform layer for iOS (POSIX files and \code{NSURLSession}); they can only be
  built on macOS;
\item \code{linuxX64()}, built only with the Gradle option \code{-Pwfdbx.linux=true},
  with the Kotlin/Native toolchain.
\end{itemize}
The native code (iOS and Linux) is not compiled by the tests of the repository, which
build the JVM code and check that the common code, which holds all of the library's
logic, uses only common APIs.

\subsection{Building and installing the library}

The library is built from a copy of the repository and installed in the local Maven
repository (\file{\textasciitilde/.m2}). In the directory \file{kotlin}:

\begin{lstlisting}[style=shell]
./gradlew publishToMavenLocal
\end{lstlisting}

A Gradle project (an Android app, a JVM program, or the common source set of a Kotlin
Multiplatform project) that lists \code{mavenLocal()} among its repositories then adds it
to its dependencies:

\begin{lstlisting}[language=Kotlin]
dependencies {
    implementation("io.github.alphanumericslab:wfdbx:1.2.0")
}
\end{lstlisting}

The build downloads the Kotlin Gradle plugin from Maven Central. Where Maven Central
cannot be reached, the Gradle property \code{mavenMirror} names another Maven repository,
for plugins and dependencies alike, on the command line or in
\file{\textasciitilde/.gradle/gradle.properties}:

\begin{lstlisting}[style=shell]
./gradlew -PmavenMirror=https://maven-central.storage-download.googleapis.com/maven2/ publishToMavenLocal
\end{lstlisting}

The Kotlin package is \code{wfdbx} (\code{import wfdbx.*}); the C standard library
emulation is in \code{wfdbx.cstd}, the tools in \code{wfdbx.apps}, and the command-line
driver in \code{wfdbx.cli} (\cref{tab:kotlin-packages}).

\subsection{The command-line tools}

The tools need the runnable jar \file{kotlin/build/libs/wfdbx-all.jar}, built in the
directory \file{kotlin} by:

\begin{lstlisting}[style=shell]
./gradlew fatJar
bin/wfdbx --version                # WFDB-x 1.2.0
bin/wfdbx --list                   # the tools
\end{lstlisting}

\subsection{The sample record}
\label{sec:kotlin-resources}

The repository includes the sample record \code{100s} in \file{tests/data}: the first
minute of record 100 of the MIT-BIH Arrhythmia Database, two ECG signals (\code{MLII} and
\code{V5}) sampled at 360\,Hz, with the reference beat annotations (annotator
\code{atr}). The library jar does not include it: an app copies the files it needs
(\file{100s.hea}, \file{100s.dat}, \file{100s.atr}) to its own storage. Work that uses the
MIT-BIH Arrhythmia Database cites it as described on PhysioNet.

\subsection{Android and iOS}

On Android and iOS, the current directory is not a useful place for files: the input and
output directories are set to a directory of the app, such as
\code{context.filesDir.path} on Android (\cref{sec:kotlin-settings}). Reading URLs on
Android needs the \code{INTERNET} permission; files are downloaded with
\code{java.net.HttpURLConnection} on the JVM and Android, and with \code{NSURLSession} on
iOS. Legacy-mode sorting of annotation files runs the \code{sortann} program, which an
iOS app cannot do (the sort of normal mode, done in the library, works everywhere).

% ==========================================================================================
\section{Quick start}
\label{sec:kotlin-quickstart}

This example reads the header, ten seconds of the signals and the annotations of the
sample record \code{100s}, and computes the mean heart rate. It then writes the first ten
seconds as a new record, \code{first10}, with a metadata entry in its header. It assumes
that the input and output directories are a copy of \file{tests/data}
(\code{wfdbxSetInputDir("tests/data")}, as in \file{kotlin/README.md}, reads the sample
record from the repository).

\begin{lstlisting}[language=Kotlin,style=example]
import wfdbx.*

// The header only: no samples are read.
val hdr = readHeader("100s")
println("${hdr.nsig} ${hdr.fs} ${hdr.nsamples}")       // 2 360.0 21600

// Seconds 10 to 20 of both signals, in ADC and physical units, frame by frame.
val rec = readRecord("100s", from = TimeSpec("0:10"), to = TimeSpec("0:20"), physical = true)
println("${rec.names} ${rec.nsamples}")                // [MLII, V5] 3600
println("${rec.physical!![0]} ${rec.physicalAt(0, 1)}")    // -0.39 -0.275

// The beat annotations, and the mean heart rate.
val ann = readAnnotations("100s", "atr")
println("${ann.size} ${ann.symbols.take(3)} ${ann.samples.take(3)}")   // 75 [+, N, N] [18, 77, 370]
val beats = (0 until ann.size).filter { isqrs(ann.codes[it]) != 0 }.map { ann.samples[it] }
val rr = (1 until beats.size).map { (beats[it] - beats[it - 1]) / ann.fs }
println("${beats.size} beats, ${"%.1f".format(60 / rr.average())} beats per minute")   // 74 beats, 73.9 ...

// Write the first 10 seconds, in ADC units, as the record "first10".
val x = readRecord("100s", to = TimeSpec("0:10"))
writeRecord("first10", x.samples, x.nsig, x.fs, fmt = 16, gains = x.gains.toDoubleArray(),
            baselines = x.baselines.toIntArray(), adczeros = x.siginfo.map { it.adczero }.toIntArray(),
            adcres = x.siginfo.map { it.adcres }.toIntArray(), units = listOf("mV"), names = x.names,
            metadata = listOf("Source" to "MIT-BIH record 100"))

// Header metadata: "# Key: value" lines of the header.
setMetadata("first10", "Age", "69")
println(readMetadata("first10"))     // [(Source, MIT-BIH record 100), (Age, 69)]
println(checkRecord("first10"))      // no problems found
\end{lstlisting}

The header \file{100s.hea} gives no units: the example writes \code{"mV"} explicitly,
since WFDB-x never assumes units, a gain, an ADC resolution or a sampling frequency that
a record does not give (\cref{chap:behavior}).

% ==========================================================================================
\section{The high-level API}
\label{sec:kotlin-api}

The high-level API is made of the functions of this section and of the classes of their
results (\cref{sec:kotlin-results}), all in the package \code{wfdbx}. They are defined in
the files \file{Record.kt}, \file{Check.kt} and \file{Metadata.kt} of
\file{kotlin/src/commonMain/kotlin/wfdbx}, whose KDoc comments document each.

\subsection{Conventions}
\label{sec:kotlin-conventions}

\paragraph{The library state: \code{Wfdb} objects and the default instance.} The state of
the library (the open record, the annotators, the input and output directories, network
access, legacy mode) is held by an object of the class \code{Wfdb}. Each high-level
function is an extension function of \code{Wfdb}, such as
\code{fun Wfdb.readRecord(...)}, called on an object: \code{w.readRecord("100s")}. Each is
also a top-level function with the same name and the same parameters, such as
\code{readRecord("100s")}, which runs it on the default instance of the calling thread,
\code{wfdbDefault()} (one per thread; \code{withDefault \{ w -> ... \}} runs a block with
it). The entries of this section give the declarations of the extension functions; the
top-level functions have the same parameters, defaults and results. Separate \code{Wfdb}
objects share no state (except the process-wide list of open output files and the cache
of fetched URLs), so several records can be read at once, in one thread or in several; a
\code{Wfdb} object is used by one thread at a time.

\begin{lstlisting}[language=Kotlin]
val w = Wfdb()                          // a library state of its own
w.wfdbxSetInputDir("tests/data")
val a = w.readRecord("100s")            // on w
val b = readRecord("100s", inputDir = "tests/data")   // on the default instance
\end{lstlisting}

\paragraph{Named arguments.} The options are parameters with default values, given as
named arguments (\code{physical = true}); there are no option structures. A \code{null}
default means ``not given'' (for \code{inputDir} and \code{network}: the current
setting).

\paragraph{Record names and directories.} A record is named as in the other parts of
WFDB-x (\cref{chap:records,chap:files}), without the \file{.hea} extension (the readers,
and \code{writeAnnotations}, also accept a trailing \code{.hea}; \code{writeRecord} does
not). A record name is a full path, absolute (\code{"/data/mitdb/100"}) or in the home
directory (\code{"\textasciitilde/Documents/100s"}), used as given; a relative path
(\code{"100s"}, \code{"data/100s"}), read from the input directory if one is set,
otherwise from the current directory, and written to the output directory if one is set,
otherwise to the current directory (missing directories are created); or a URL
(\code{"https://physionet.org/files/mitdb/1.0.0/100"}), read only if network access is
on, which cannot be written. A writer given a name with a path creates the files in that
directory, and the header names the record and its signal file without the path
(\code{rec}, \code{rec.dat}), so that the record opens under any path that leads to it;
in legacy mode, as in the original, the header gives the name with its path. The input
and output directories are the current directory unless they are set, with \code{wfdbxSetInputDir} and \code{wfdbxSetOutputDir}
(\cref{sec:kotlin-settings}) or with the environment variables \code{WFDBX\_INPUT\_DIR}
and \code{WFDBX\_OUTPUT\_DIR}. URLs are read only when network access is on. The
functions that read take the parameters \code{inputDir} and \code{network}, and those
that write take \code{dir}: they apply to that one call, and the previous settings are
restored when it returns, also after an error. \code{inputDir} is optional, and applies
to relative names only, not to full paths or URLs; with \code{dir}, the record name must
be a plain name (\code{"rec"}), and the record opens from that directory under it. The signal files that a header names are
read from the header's directory. Record, annotator and directory names are converted to
bytes with the platform's file name encoding.

\paragraph{Times.} The parameters \code{from} and \code{to} are \code{TimeSpec} values: a
sample number (a \code{TimeSpec.Sample}, made from an \code{Int} or a \code{Long}), or a
WFDB time string, as \code{strtim} reads it (a \code{TimeSpec.Text}):
\begin{lstlisting}[language=Kotlin]
TimeSpec(3600)            // sample 3600
TimeSpec("0:10")          // 10 seconds; also "1:02:03.5"
TimeSpec("s3600")         // sample 3600
TimeSpec("e")             // the end of the record
TimeSpec("[14:00:00]")    // a time of day
\end{lstlisting}
A time in seconds needs the sampling frequency; when the record does not give it, such a
time throws a \code{WfdbException}.

\paragraph{Samples, frame by frame.} The samples of a record are one array, frame by
frame, as the library returns them: \code{samples[i * nsig + s]} is sample \code{i} of
signal \code{s} (\code{rec.at(i, s)}), so that a record of millions of samples is one
array, without a copy per signal. A signal with \code{spf} samples per frame has
\code{spf} consecutive values in each frame.

\paragraph{Per-signal and per-annotation arguments.} The arguments that describe each
signal (\code{gains}, \code{units}, \ldots) take one value for all the signals, or exactly
one value per signal. Those that describe each annotation (\code{subtypes}, \code{chans},
\ldots) have exactly one value per annotation.

\paragraph{Text.} Text in the files (signal descriptions, units, info strings, metadata,
aux strings, mnemonics) is a \code{String} with one character per byte (ISO-8859-1), so
that every byte round-trips. The fields of \code{Siginfo} are byte strings
(\code{ByteArray}), as in C; \code{latin1()}, from \code{wfdbx.cstd}, converts a
\code{ByteArray} to a \code{String} and back.

\paragraph{Signal parameters.} No sampling frequency, gain, ADC resolution or units is
ever assumed. Where one is needed and the record does not give it, it must be given as an
argument; otherwise the function throws a \code{WfdbException}. A record whose header
gives no sampling frequency has \code{fs = 0}.

\paragraph{Legacy mode.} With \code{legacy = true}, a function behaves as in the original
WFDB package, with its assumed signal parameters and its search for files
(\cref{chap:behavior}); legacy mode is on, on that \code{Wfdb} object, for the duration
of the call only. With \code{legacy = false} (the default), the current setting of
\code{wfdbxSetLegacy} applies, which is off unless it was turned on.

\paragraph{State.} The functions that read and write records and annotation files close
anything opened with the WFDB library functions on the same \code{Wfdb} object (they call
\code{wfdbquit()}) before they start, and close their own files before they return,
whatever the outcome. They leave the time resolution of output annotators
(\code{getafreq()}) and the getvec mode (\code{getgvmode()}) as they were, also after an
error and after reading an EDF record, whose opening sets both, as in C (\code{wfdbquit}
resets neither): \code{readRecord} sets the getvec mode for its own reading only, and the
\code{fs} of \code{writeAnnotations} applies to that file only. The metadata functions
open no signals and leave the library's state as it is.

\paragraph{The examples} of this section use the top-level functions, with the input and
output directories set to a copy of \file{tests/data} (\code{wfdbxSetInputDir(d)} and
\code{wfdbxSetOutputDir(d)}), which holds the test records of WFDB-x, among them
\code{100s}; they write records of their own there. The listings of style \code{example}
were compiled and run with the library on the JVM, and print what their comments say.

% ------------------------------------------------------------------------------------------
\subsection{Reading records}
\label{sec:kotlin-reading}

\apientry{readRecord}{Read the signals of a record}
\begin{lstlisting}[language=Kotlin,style=signature]
fun Wfdb.readRecord(
    record: String,
    channels: List<String>? = null,
    from: TimeSpec? = null,
    to: TimeSpec? = null,
    physical: Boolean = false,
    highres: Boolean = false,
    pad: Boolean = false,
    checksumWarning: Boolean = true,
    onWarning: ((String) -> Unit)? = null,
    inputDir: String? = null,
    network: Boolean? = null,
    legacy: Boolean = false,
): Record
\end{lstlisting}
Reads the signals of a record, or part of it, into one \code{IntArray} of ADC units,
frame by frame, together with the description of the record. The samples are those that
\code{getvec} returns (\cref{chap:library}), read in large blocks. Invalid samples are
\code{WFDB\_INVALID\_SAMPLE} ($-32768$) in ADC units, and NaN in physical units. When the
read reaches the end of the record, the samples are compared with the checksums of the
header (\cref{sec:kotlin-warnings}). On any error once the record is open, its input
signals are closed, and the getvec mode and the settings changed for the call are
restored.
\begin{params}
\item[record] Record name, without the \file{.hea} extension (a trailing \code{".hea"} is
  also accepted): a full path, absolute (\code{"/data/mitdb/100"}) or in the home
  directory (\code{"\textasciitilde/Documents/100s"}), used as given; a relative path
  (\code{"100s"}, \code{"data/100s"}), read from the input directory if one is set,
  otherwise from the current directory; or a URL
  (\code{"https://physionet.org/files/mitdb/1.0.0/100"}), read only if network access is
  on.
\item[channels] The signals to read, in the order in which they are returned: a list of
  signal numbers (\code{"0"}, \code{"1"}, \ldots) or signal descriptions (as in
  \code{Record.names}, such as \code{"MLII"}), which may be mixed. A name is looked up
  among the descriptions first, then with \code{findsig}. Default \code{null}: all the
  signals, in the order of the header.
\item[from] The first sample to read: a \code{TimeSpec}, a sample number ($\ge 0$) or a
  time string. Default \code{null}: the beginning of the record.
\item[to] The sample after the last one to read: a \code{TimeSpec}. Default \code{null}:
  the end of the record.
\item[physical] \code{true} to return the samples in physical units too, in
  \code{Record.physical} (a \code{DoubleArray} in the layout of \code{samples}), computed
  as $(\mathit{adc} - \mathit{baseline}) / \mathit{gain}$. Every signal read must then have
  a gain. Default \code{false}.
\item[highres] \code{false} (the default) to read a multi-frequency record at its frame
  frequency, with the samples of the faster signals averaged to one per frame;
  \code{true} to read every sample, at the highest sampling frequency of the record (the
  samples of the slower signals are repeated). \code{Record.fs} is then that frequency.
\item[pad] \code{true} to replace each invalid sample by the previous valid sample of the
  same signal. Default \code{false}.
\item[checksumWarning] \code{true} (the default) to warn when the samples do not match
  the checksums of the header; \code{false} for no warning. The samples are returned
  either way, and \code{checksumOk} of the result is \code{false}.
\item[onWarning] A function that receives the warning (without the \code{WFDB-x: }
  prefix and the newline). Default \code{null}: the library prints \code{WFDB-x: }
  followed by the warning on the standard error, unless \code{wfdbquiet()} was called.
\item[inputDir] The input directory for this call only (optional): a local directory, a
  \code{\textasciitilde} path or a URL, under which relative record names are read. It
  does not apply to full paths or URLs. Default \code{null}: the current setting
  (\code{wfdbxSetInputDir}).
\item[network] \code{true} to allow reading URLs for this call only, \code{false} to
  forbid it. Default \code{null}: the current setting (\code{wfdbxSetNetwork}).
\item[legacy] \code{true} to behave as the original WFDB package (legacy mode, for this
  call only). Default \code{false}.
\end{params}
\returns{A \code{Record} (\cref{tab:kotlin-record}).}
\raises{\code{WfdbException} if the header or a signal file cannot be read (``can't read
record REC''), if a channel is not in the record (``record REC has no signal C''), if
\code{from} is a negative sample number (``readRecord: from must not be negative''), if a
time string needs the sampling frequency and the header gives none (``sampling frequency
of record REC is not specified''), if \code{physical} is \code{true} and a signal read has
no gain (``signal N of record REC is uncalibrated''), if \code{from} is beyond the end of
the record (``can't seek to sample N of record REC''), or if a signal file is shorter than
its header says (``can't read signal data of record REC (getvec status -3)''). Warns on a
checksum error (see \code{checksumWarning} and \code{onWarning}).}
\begin{lstlisting}[language=Kotlin,style=example]
import wfdbx.*

// Seconds 10 to 20 of lead MLII, in physical units.
var rec = readRecord("100s", channels = listOf("MLII"), from = TimeSpec("0:10"),
                     to = TimeSpec("0:20"), physical = true)
println("${rec.nsamples} ${rec.nsig} ${rec.names} ${rec.fs} ${rec.sampfrom}")   // 3600 1 [MLII] 360.0 3600

// The first 3 samples of both signals, in ADC units, V5 first.
rec = readRecord("100s", channels = listOf("1", "0"), to = TimeSpec(3))
println(rec.samples.toList())       // [1011, 995, 1011, 995, 1011, 995]
println("${rec.names} ${rec.gains} ${rec.baselines} ${rec.units}")
// [V5, MLII] [200.0, 200.0] [1024, 1024] [null, null]

// The whole record: the checksums are verified.
rec = readRecord("100s")
println("${rec.nsamples} ${rec.checksumOk}")      // 21600 true
\end{lstlisting}
A multi-frequency record (here the test record \code{mfreq}, whose signals have 1, 2
and 4 samples per frame at a frame frequency of 100\,Hz) is read at its frame frequency,
or at its highest sampling frequency:
\begin{lstlisting}[language=Kotlin,style=example]
import wfdbx.*

var rec = readRecord("mfreq")
println("${rec.fs} ${rec.nsamples} ${rec.nsig}")     // 100.0 600 3
rec = readRecord("mfreq", highres = true)
println("${rec.fs} ${rec.nsamples} ${rec.nsig}")     // 400.0 2400 3
println(rec.siginfo.map { it.spf })                  // [1, 2, 4]
\end{lstlisting}

\apientry{readHeader}{Read the header of a record}
\begin{lstlisting}[language=Kotlin,style=signature]
fun Wfdb.readHeader(
    record: String,
    inputDir: String? = null,
    network: Boolean? = null,
    legacy: Boolean = false,
    onWarning: ((String) -> Unit)? = null,
): Header
\end{lstlisting}
Reads the header of a record, without opening its signal files. The result describes the
record and each of its signals, and holds the info strings and the metadata entries of the
header (of the \file{.hea} file only, not of a \file{.info} file; none for an EDF file or a
header read from the standard input). For a multi-segment record, the master header is
read, and \code{nsamples} is the length of the whole record.
\begin{params}
\item[record] Record name, without the \file{.hea} extension (a trailing \code{".hea"} is
  also accepted): a full path, absolute (\code{"/data/mitdb/100"}) or in the home
  directory (\code{"\textasciitilde/Documents/100s"}), used as given; a relative path
  (\code{"100s"}, \code{"data/100s"}), read from the input directory if one is set,
  otherwise from the current directory; or a URL
  (\code{"https://physionet.org/files/mitdb/1.0.0/100"}), read only if network access is
  on.
\item[inputDir] The input directory for this call only (optional): a local directory, a
  \code{\textasciitilde} path or a URL, under which relative record names are read. It
  does not apply to full paths or URLs. Default \code{null}: the current setting
  (\code{wfdbxSetInputDir}).
\item[network] \code{true} to allow reading URLs for this call only, \code{false} to
  forbid it. Default \code{null}: the current setting.
\item[legacy] \code{true} to read the header as the original WFDB package does. Default
  \code{false}.
\item[onWarning] A function that receives the warnings about \code{Header.metadata} (a
  key given more than once, \cref{sec:kotlin-warnings}). Default \code{null}: the library
  prints \code{WFDB-x: } followed by the warning on the standard error, unless
  \code{wfdbquiet()} was called.
\end{params}
\returns{A \code{Header} (\cref{tab:kotlin-header}).}
\raises{\code{WfdbException} if the header cannot be found or read (``can't read header
for record REC'').}
\begin{lstlisting}[language=Kotlin,style=example]
import wfdbx.*
import wfdbx.cstd.latin1

val hdr = readHeader("100s")
println("${hdr.name} ${hdr.fs} ${hdr.nsig} ${hdr.nsamples}")    // 100s 360.0 2 21600
for (s in hdr.siginfo) {
    println("${s.desc?.latin1()} ${s.fname?.latin1()} ${s.fmt} ${s.gain} ${s.adcres} ${s.units?.latin1()}")
}
// MLII 100s.dat 212 200.0 11 null
// V5 100s.dat 212 200.0 11 null
println("\"${hdr.comments[0]}\"")   // " 69 M 1085 1629 x1": the first info string
println(hdr.baseTime)               // null: the header gives no start time
\end{lstlisting}

\apientry{readAnnotations}{Read an annotation file}
\begin{lstlisting}[language=Kotlin,style=signature]
fun Wfdb.readAnnotations(
    record: String,
    annotator: String,
    from: TimeSpec? = null,
    to: TimeSpec? = null,
    fs: Double? = null,
    inputDir: String? = null,
    network: Boolean? = null,
    legacy: Boolean = false,
): Annotations
\end{lstlisting}
Reads the annotations of an annotation file, in file order. The sampling frequency,
needed to convert times in seconds, is that of the record's header, or, for a record
without a header or whose header gives none, the argument \code{fs}.
\begin{params}
\item[record] Record name, without the \file{.hea} extension (a trailing \code{".hea"} is
  also accepted): a full path, absolute (\code{"/data/mitdb/100"}) or in the home
  directory (\code{"\textasciitilde/Documents/100s"}), used as given; a relative path
  (\code{"100s"}, \code{"data/100s"}), read from the input directory if one is set,
  otherwise from the current directory; or a URL
  (\code{"https://physionet.org/files/mitdb/1.0.0/100"}), read only if network access is
  on. Its header, if there is one, gives the sampling frequency.
\item[annotator] Annotator name: the suffix of the annotation file. \code{"atr"} reads the
  file \file{RECORD.atr}.
\item[from] The first time to read: a \code{TimeSpec}. Default \code{null}: the beginning
  of the file.
\item[to] The time at which to stop: annotations at or after it are not read. Default
  \code{null}: the end of the file.
\item[fs] The sampling frequency, in Hz ($> 0$), to use when the record has no header or
  its header gives none; it is ignored otherwise. Default \code{null}: none.
\item[inputDir] The input directory for this call only (optional): a local directory, a
  \code{\textasciitilde} path or a URL, under which relative record names are read. It
  does not apply to full paths or URLs. Default \code{null}: the current setting
  (\code{wfdbxSetInputDir}).
\item[network] \code{true} to allow reading URLs for this call only, \code{false} to
  forbid it. Default \code{null}: the current setting.
\item[legacy] \code{true} to behave as the original WFDB package (which assumes 250\,Hz
  for a record without a frequency). Default \code{false}.
\end{params}
\returns{An \code{Annotations} object (\cref{tab:kotlin-annotations}). Its \code{fs} is
the frequency in which the times are counted: the one written in the annotation file, if
there is one, or else the sampling frequency; it is 0 if none of them is known.}
\raises{\code{WfdbException} if \code{fs} is given and is not $> 0$ (``readAnnotations:
fs must be > 0''), if the annotation file cannot be opened (``can't read annotator ANN for
record REC''), if a time string needs the sampling frequency and none is known, or if the
file ends before its end-of-file marker (``can't read annotation file ANN of record REC
(getann status -3)'').}
\begin{lstlisting}[language=Kotlin,style=example]
import wfdbx.*

var ann = readAnnotations("100s", "atr", from = TimeSpec("0:10"), to = TimeSpec("0:30"))
println("${ann.size} ${ann.fs}")                              // 24 360.0
println("${ann.samples.take(3)} ${ann.symbols.take(3)}")      // [3862, 4170, 4466] [N, N, N]

// The whole file: the first annotation is a rhythm change, with an aux string.
ann = readAnnotations("100s", "atr")
println("${ann.symbols[0]} ${ann.codes[0] == RHYTHM} ${ann.aux[0]}")   // + true (N
println(ann.symbols.toSortedSet())                            // [+, A, N]
\end{lstlisting}
For a record without a header, \code{fs} gives the frequency in which times in seconds
are converted (the library reports on the standard error that it cannot open the
header):
\begin{lstlisting}[language=Kotlin,style=example]
import wfdbx.*

writeAnnotations("beats", "qrs", longArrayOf(100, 460, 820), symbols = listOf("N", "N", "N"))
val ann = readAnnotations("beats", "qrs", to = TimeSpec("0:02"), fs = 360.0)
println("${ann.samples.toList()} ${ann.fs}")        // [100, 460] 360.0
\end{lstlisting}

% ------------------------------------------------------------------------------------------
\subsection{Writing records}
\label{sec:kotlin-writing}

\apientry{writeRecord}{Write a header and a signal file}
\begin{lstlisting}[language=Kotlin,style=signature]
fun Wfdb.writeRecord(
    record: String,
    samples: IntArray,
    nsig: Int,
    fs: Double,
    fmt: Int = 16,
    gains: DoubleArray? = null,
    baselines: IntArray? = null,
    adczeros: IntArray? = null,
    adcres: IntArray? = null,
    units: List<String?>? = null,
    names: List<String>? = null,
    baseTime: String? = null,
    comments: List<String>? = null,
    sampsPerFrame: IntArray? = null,
    dir: String? = null,
    legacy: Boolean = false,
    metadata: List<Pair<String, String?>>? = null,
)

fun Wfdb.writeRecord(
    record: String,
    samples: DoubleArray,
    nsig: Int,
    fs: Double,
    fmt: Int = 16,
    gains: DoubleArray? = null,
    baselines: IntArray? = null,
    adczeros: IntArray? = null,
    adcres: IntArray? = null,
    units: List<String?>? = null,
    names: List<String>? = null,
    baseTime: String? = null,
    comments: List<String>? = null,
    sampsPerFrame: IntArray? = null,
    dir: String? = null,
    legacy: Boolean = false,
    metadata: List<Pair<String, String?>>? = null,
)
\end{lstlisting}
Writes samples as a new record: the header \file{RECORD.hea} and one signal file,
\file{RECORD.dat}, which holds all the signals. Existing files of these names are
overwritten. The first overload takes the samples in ADC units (an \code{IntArray}), the
second in physical units (a \code{DoubleArray}), converted to ADC units as
$\mathrm{round}(x \cdot \mathit{gain}) + \mathit{baseline}$, with halves rounded away from
zero; NaN becomes an invalid sample, and an infinite value, or one whose ADC value is
beyond the range of 32-bit integers, is an error. The sampling frequency, the gain and
ADC resolution of each signal, and the units of each calibrated signal must be given:
there are no defaults. The per-signal arguments (\code{gains}, \code{baselines},
\code{adczeros}, \code{adcres}, \code{units}, \code{names}, \code{sampsPerFrame}) take one
value for all the signals, or one value per signal. All the arguments are checked before
any file is created, and the output signals are closed whatever happens.
\begin{params}
\item[record] Record name (not empty), without the \file{.hea} extension: a plain name
  (\code{"rec"}) or a relative path (\code{"data/rec"}), created in the output directory
  if one is set, otherwise in the current directory, missing directories included; or a
  full path, absolute or in the home directory (\code{"\textasciitilde/Documents/rec"}),
  used as given. With a path, the files are created in that directory, and the header
  names the record and its signal file without the path (\code{rec}, \code{rec.dat}), so
  that the record opens under any path that leads to it (in legacy mode, as in the
  original, the header gives the name with its path). A URL cannot be written.
\item[samples] The samples, frame by frame: signal 0, signal 1, \ldots{} of frame 0, then
  those of frame 1, and so on (a signal with \code{spf} samples per frame has \code{spf}
  consecutive values in each frame). An \code{IntArray} of ADC units, or a
  \code{DoubleArray} of physical units. Its size must be a multiple of the frame size.
\item[nsig] The number of signals (at least 1).
\item[fs] The sampling frequency, in Hz (frames per second): a finite number $> 0$.
\item[fmt] The storage format of all the signals, such as \code{16}, \code{212},
  \code{24} or \code{516} (FLAC) (\cref{chap:formats}). Default \code{16}.
\item[gains] The gain of each signal, in ADC units per physical unit: a finite number;
  \code{0} marks an uncalibrated signal, which is allowed only for samples given in ADC
  units. Negative gains are allowed. Required (default \code{null}: not given).
\item[baselines] The ADC value that corresponds to 0 physical units. Default
  \code{null}: the ADC zero of the signal.
\item[adczeros] The ADC value at the middle of the ADC range. Default \code{null}: 0.
\item[adcres] The ADC resolution, in bits ($\ge 1$). Required (default \code{null}: not
  given).
\item[units] The physical units, such as \code{"mV"}; required for each signal whose gain
  is not 0 (\code{null} for an uncalibrated signal). Default \code{null}: not given.
\item[names] The signal descriptions, such as \code{"MLII"}. Default \code{null}: no
  descriptions.
\item[baseTime] The start time, with an optional date: \code{"hh:mm:ss"} or
  \code{"hh:mm:ss dd/mm/yyyy"}. Default \code{null}: no start time.
\item[comments] Info strings, written at the end of the header as comment lines
  (\code{"\#"} followed by the string). Default \code{null}: none.
\item[sampsPerFrame] For a multi-frequency record, the number of samples per frame of
  each signal (integers $\ge 1$); \code{fs} is then the frame frequency. Default
  \code{null}: one sample per frame for every signal.
\item[dir] A directory (which must exist) in which to create the files, for this call
  only. \code{record} must then be a plain name, without a directory, and the record opens
  from that directory under its plain name. Default \code{null} (or \code{""}): the output
  directory.
\item[legacy] \code{true} to behave as the original WFDB package: the gain defaults to
  200, the ADC resolution to one based on the format, and no units are required. Default
  \code{false}.
\item[metadata] Metadata entries (\cref{chap:metadata}) to write to the header after the
  comments, as \code{\# Key: value} (through \code{putinfo}), in the order given: a list
  of pairs, \code{listOf("Age" to "63")} (a map converts with \code{toList()}). Keys and
  values follow the rules of \code{setMetadata}; no two keys may be equal (ignoring
  letter case), none may be the key of an entry in \code{comments}, and a \code{null}
  value is not allowed. Default \code{null}: none.
\end{params}
\returns{Nothing (\code{Unit}).}
\raises{\code{WfdbException}, before any file is created: ``writeRecord: sampsPerFrame
must hold integers >= 1'' (checked first, also in legacy mode), ``writeRecord: the record
name must not be empty'', ``writeRecord: at least one signal is required'',
``writeRecord: NAME must have one value per signal'', ``writeRecord: the sampling
frequency must be > 0'' (not checked in legacy mode), ``writeRecord: the sampling
frequency must be finite'', ``writeRecord: gains must be specified'', ``writeRecord:
gains must be finite'', ``writeRecord: ADC resolution(s) must be specified'',
``writeRecord: units must be specified'', ``writeRecord: incorrect base time 'TIME'{}'',
``writeRecord: the number of samples is not a multiple of the frame size'',
``writeRecord: signal N is uncalibrated'' and ``writeRecord: sample I is infinite or
beyond the range of 32-bit ADC units'' (physical units), ``record name must not contain a
directory when dir is given'', ``directory DIR does not exist''; a
\code{WfdbArgumentException} if an entry of \code{metadata} is not valid (``writeRecord:
metadata key 'a:b' must not contain ':'{}'', ``writeRecord: value of metadata key 'Age'
must not be null'', ``writeRecord: metadata key 'Age' is also given in comments'',
\ldots). \code{WfdbException} if the signal file or the header cannot be created
(``can't create signal file for record REC'', ``can't create header for record REC''),
or cannot be written completely, for example on a full disk (``write error on record
REC'').}
\begin{lstlisting}[language=Kotlin,style=example]
import wfdbx.*
import kotlin.math.abs

// Two signals of random numbers, in millivolts, at 360 Hz: the record "rand"
// (rand.hea and rand.dat).
val rnd = kotlin.random.Random(1)
val x = DoubleArray(3600 * 2) { rnd.nextDouble(-3.0, 3.0) }
writeRecord("rand", x, nsig = 2, fs = 360.0, fmt = 16, gains = doubleArrayOf(1000.0),
            adcres = intArrayOf(16), units = listOf("mV"), names = listOf("A", "B"),
            baseTime = "14:00:03 21/06/2024", comments = listOf("random numbers"),
            metadata = listOf("Generator" to "kotlin.random"))
val rec = readRecord("rand", physical = true)
println("${rec.nsamples} ${rec.baseTime}")     // 3600 [14:00:03 21/06/2024]
println(x.indices.all { abs(rec.physical!![it] - x[it]) <= 0.0005 })   // true: rounded to 1/1000 mV
println(readHeader("rand").comments)           // [random numbers,  Generator: kotlin.random]

// An uncalibrated signal (gain 0, no units), in ADC units, in the directory "out".
writeRecord("out/counts", IntArray(1000) { it }, nsig = 1, fs = 100.0, gains = doubleArrayOf(0.0),
            adcres = intArrayOf(12))
println(readRecord("out/counts").samples.take(3))    // [0, 1, 2]
\end{lstlisting}
A multi-frequency record: signal \code{fast} has 2 samples per frame, signal
\code{slow} one; each frame holds the two samples of \code{fast} and then the sample of
\code{slow}:
\begin{lstlisting}[language=Kotlin,style=example]
import wfdbx.*

val frames = 100                                    // 1 second at 100 frames per second
val x = IntArray(3 * frames) { if (it % 3 == 2) 500 else 2 * (it / 3) + it % 3 }
writeRecord("mf", x, nsig = 2, fs = 100.0, gains = doubleArrayOf(200.0), adcres = intArrayOf(12),
            units = listOf("mV"), names = listOf("fast", "slow"), sampsPerFrame = intArrayOf(2, 1))
val rec = readRecord("mf", highres = true)
println("${rec.fs} ${rec.samples.take(6)}")        // 200.0 [0, 500, 1, 500, 2, 500]
\end{lstlisting}

\apientry{writeAnnotations}{Write an annotation file}
\begin{lstlisting}[language=Kotlin,style=signature]
fun Wfdb.writeAnnotations(
    record: String,
    annotator: String,
    samples: LongArray,
    symbols: List<String>? = null,
    codes: IntArray? = null,
    subtypes: IntArray? = null,
    chans: IntArray? = null,
    nums: IntArray? = null,
    aux: List<String?>? = null,
    fs: Double? = null,
    sort: Boolean? = null,
    dir: String? = null,
    legacy: Boolean = false,
)

fun Wfdb.writeAnnotations(
    record: String,
    annotator: String,
    a: Annotations,
    fs: Double? = null,
    sort: Boolean? = null,
    dir: String? = null,
    legacy: Boolean = false,
)
\end{lstlisting}
Writes an annotation file, \file{RECORD.ANNOTATOR}, in the MIT format; an existing file of
that name is overwritten. The record need not have a header. The annotation types are
given as \code{symbols} or as \code{codes} (\cref{chap:formats}). If the annotations are
not in time order, the file is sorted when it is written, unless \code{sort = false}. The
second overload writes the annotations of an \code{Annotations} object, as
\code{readAnnotations} returns them: its samples, codes, subtypes, channels, annotator
numbers and aux strings (its record, annotator and \code{fs} are not used). All the
arguments are checked before the file is created.
\begin{params}
\item[record] Record name, without the \file{.hea} extension (a trailing \code{".hea"} is
  also accepted): a plain name or a relative path (\code{"data/100s"}) is created in the
  output directory if one is set, otherwise in the current directory; a full path,
  absolute or in the home directory (\code{"\textasciitilde/Documents/100s"}), is used as
  given. The file is \file{RECORD.ANNOTATOR} at that place (\file{data/100s.qrs}). A URL
  cannot be written.
\item[annotator] Annotator name: the suffix of the file, made of letters, digits,
  \code{\_}, \code{-} and \code{\textasciitilde}.
\item[samples] The annotation times, in samples, one per annotation.
\item[symbols] The mnemonics of the annotation types, such as \code{"N"} or \code{"V"},
  one per annotation. An unknown mnemonic becomes code 0 (\code{NOTQRS}). Default
  \code{null}: the types are given by \code{codes}.
\item[codes] The annotation codes (\code{NORMAL}, \code{PVC}, \ldots), one per annotation,
  instead of \code{symbols}. If both are given (as \code{readAnnotations} returns them),
  each symbol must be that of its code (an empty symbol is not checked). Default
  \code{null}.
\item[subtypes] The subtype of each annotation, stored as a signed byte (other values wrap
  around to $-128 \ldots 127$, as in C). Default \code{null}: 0.
\item[chans] The channel number of each annotation, stored as an unsigned byte (other
  values wrap around to $0 \ldots 255$). Default \code{null}: 0.
\item[nums] The annotator number of each annotation, stored as a signed byte. Default
  \code{null}: 0.
\item[aux] The auxiliary text of each annotation, with \code{null} where there is none.
  Texts longer than 255 bytes are cut. Default \code{null}: none.
\item[fs] If given (a finite number $> 0$), the frequency in which the times are counted,
  written to the file. Default \code{null}: none is written.
\item[sort] \code{true} to always sort the file, \code{false} to keep the given order
  (without the library's advice to run \code{sortann}), for this call only. Default
  \code{null}: the environment variable \code{WFDBANNSORT} decides (sorting is on unless
  it is \code{0}).
\item[dir] A directory (which must exist) in which to create the file, for this call
  only. \code{record} must then be a plain name, without a directory: the file is
  \file{DIR/RECORD.ANNOTATOR}. Default \code{null}: the output directory.
\item[legacy] \code{true} to behave as the original WFDB package. Default \code{false}.
\item[a] The annotations to write (second overload): an \code{Annotations} object, as
  \code{readAnnotations} returns it.
\end{params}
\returns{Nothing (\code{Unit}).}
\raises{\code{WfdbException}, before the file is created, if neither \code{symbols} nor
\code{codes} is given (``symbols or codes must be given''), if both are given and
disagree (``writeAnnotations: symbols and codes disagree at annotation K'', K counting
from 0), if an array does not have one value per annotation (``writeAnnotations: NAME
must have one value per annotation''), if \code{fs} is not finite and $> 0$
(``writeAnnotations: fs must be finite and > 0''), or as for the \code{dir} of
\code{writeRecord}; if the file cannot be created (``can't create annotator ANN for record
REC''), or cannot be written completely (``write error on annotator ANN for record
REC'').}
\begin{lstlisting}[language=Kotlin,style=example]
import wfdbx.*

// Three beats and a rhythm change, by symbol, with an aux string for the rhythm.
writeAnnotations("100s", "test", longArrayOf(18, 77, 370, 662),
                 symbols = listOf("+", "N", "V", "N"), aux = listOf("(N", null, null, null))
var ann = readAnnotations("100s", "test")
println("${ann.symbols} ${ann.aux}")       // [+, N, V, N] [(N, null, null, null]

// A copy of the reference annotations, as readAnnotations returns them.
val a = readAnnotations("100s", "atr")
writeAnnotations("100s", "copy", a)
val b = readAnnotations("100s", "copy")
println("${b.samples.contentEquals(a.samples)} ${b.symbols == a.symbols} ${b.aux == a.aux}")
// true true true

// Annotations out of time order are sorted when the file is written (the library
// reports it on the standard error).
writeAnnotations("100s", "unsorted", longArrayOf(500, 100), codes = intArrayOf(NORMAL, NORMAL))
println(readAnnotations("100s", "unsorted").samples.toList())     // [100, 500]
\end{lstlisting}

% ------------------------------------------------------------------------------------------
\subsection{Checking records}
\label{sec:kotlin-checking}

\apientry{checkRecord}{Check a record and its annotation files}
\begin{lstlisting}[language=Kotlin,style=signature]
fun Wfdb.checkRecord(
    record: String,
    annotators: List<String> = emptyList(),
    inputDir: String? = null,
    network: Boolean? = null,
): CheckReport
\end{lstlisting}
Checks the header of a record, its signal specifications, every sample of its signal
files, each segment of a multi-segment record, and the annotation files named, and
reports each problem found. A problem in the record never throws an exception: it is
reported, with a severity (\code{Severity.ERROR}: the record is damaged or does not follow
the WFDB format, and programs will read wrong data or fail; \code{Severity.WARNING}: the
record can be read, but something is missing or unusual, and may give wrong results;
\code{Severity.INFO}: a fact worth knowing, such as invalid samples). The problems, their
codes and their messages are the same in the seven languages
(\cref{tab:python-issue-codes}). The checks run in normal mode, even if legacy mode is on,
and the library prints no messages during the call. They are similar to those of Benjamin
Moody's \code{wfdbcheck}.
\begin{params}
\item[record] Record name, without the \file{.hea} extension (a trailing \code{".hea"} is
  also accepted): a full path, absolute (\code{"/data/mitdb/100"}) or in the home
  directory (\code{"\textasciitilde/Documents/100s"}), used as given; a relative path
  (\code{"100s"}, \code{"data/100s"}), read from the input directory if one is set,
  otherwise from the current directory; or a URL
  (\code{"https://physionet.org/files/mitdb/1.0.0/100"}), read only if network access is
  on (as for \code{readRecord}).
\item[annotators] The annotator names whose files are checked, such as
  \code{listOf("atr")}. Default \code{emptyList()}: none.
\item[inputDir] The input directory for this call only (optional): a local directory, a
  \code{\textasciitilde} path or a URL, under which relative record names are read. It
  does not apply to full paths or URLs. Default \code{null}: the current setting
  (\code{wfdbxSetInputDir}).
\item[network] \code{true} to allow reading URLs for this call only, \code{false} to
  forbid it. Default \code{null}: the current setting.
\end{params}
\returns{A \code{CheckReport} (\cref{tab:kotlin-checkreport}): \code{report.ok} is
\code{true} if no problem is an error, \code{report.issues} lists the problems
(\code{Issue} objects, \cref{tab:kotlin-issue}), and \code{println(report)} prints one
line per problem, or \code{no problems found}.}
\raises{\code{WfdbException} if \code{record} is empty (``record must not be empty''). A
record that cannot be read is reported (code \code{header\_unreadable}), not thrown.}
\begin{lstlisting}[language=Kotlin,style=example]
import wfdbx.*

var report = checkRecord("100s", listOf("atr"))
println(report.ok)          // true: no errors
println(report)
// warning: signal 0: the units are not specified
// warning: signal 1: the units are not specified

// A damaged record: one sample of a 12-bit signal is overwritten with 2048 (JVM file
// access, in the output directory).
writeRecord("damaged", IntArray(100) { it }, nsig = 1, fs = 100.0, gains = doubleArrayOf(100.0),
            adcres = intArrayOf(12), units = listOf("mV"), names = listOf("ramp"))
java.io.RandomAccessFile(wfdbxGetOutputDir() + "/damaged.dat", "rw").use { f ->
    f.seek(20)                                  // sample 10 (2 bytes per sample in format 16)
    f.write(byteArrayOf(0x00, 0x08))            // 0x0800 = 2048, little-endian
}
report = checkRecord("damaged")
println(report.ok)          // false
for (issue in report.issues) println("${issue.severity} ${issue.code} ${issue.where}")
// error sample_range signal 0
// error checksum signal 0
println(report.lines()[0])
// error: signal 0: 1 sample is outside the ADC range (-2048 to 2047): sample 10 (value 2048)
\end{lstlisting}

% ------------------------------------------------------------------------------------------
\subsection{Header metadata}
\label{sec:kotlin-metadata}

These functions read and change the metadata of a header: its info strings of the form
\code{\# Key: value}, such as \code{\# Age: 63} (\cref{chap:metadata}). In short: keys are
compared without regard to letter case (A--Z and a--z only); a key must not be empty and
must not contain \code{:}, \code{\#}, a line break or a NUL character, and a value must
not contain a line break or a NUL character; spaces, tabs, VT, FF and CR at either end of
keys and values are removed; text is one character per byte (ISO-8859-1), and a character
beyond U+00FF is an argument error. A change rewrites the header in one step: the new
header is written to a temporary file in the header's directory, which then replaces the
header, keeping its permission bits where the platform has them; every other byte of the
header is kept, and all the arguments are checked first, so an error leaves the header
unchanged. If nothing changes (and \code{dir} is not given), nothing is written. Only the
\file{.hea} file is read (not a \file{.info} file); for a multi-segment record, it is the
master header, unless a segment is named. Invalid arguments throw a
\code{WfdbArgumentException}, and a header that cannot be read or written a
\code{WfdbException} (\cref{sec:kotlin-warnings}). Metadata often holds information about
patients: WFDB-x stores it as plain text, exactly as given, and does not de-identify it.

\begin{phibox}
A header can hold identifying data about a patient (a name, a date of birth, a medical
record number). The metadata functions read and write such data as given. Records that
are shared must be de-identified first.
\end{phibox}

\apientry{readMetadata}{Read the metadata of a record}
\begin{lstlisting}[language=Kotlin,style=signature]
fun Wfdb.readMetadata(
    record: String,
    onWarning: ((String) -> Unit)? = null,
    inputDir: String? = null,
    network: Boolean? = null,
): List<Pair<String, String>>
\end{lstlisting}
Reads the metadata entries of the header of a record, in file order. An info string is an
entry if it contains a colon: the key is the text before the first colon, the value the
rest of the line, both without spaces at either end. Info strings without a colon are free
comments, which are not returned.
\begin{params}
\item[record] Record name, without the \file{.hea} extension (a trailing \code{".hea"} is
  also accepted): a full path, absolute (\code{"/data/mitdb/100"}) or in the home
  directory (\code{"\textasciitilde/Documents/100s"}), used as given; a relative path
  (\code{"100s"}, \code{"data/100s"}), read from the input directory if one is set,
  otherwise from the current directory; or a URL
  (\code{"https://physionet.org/files/mitdb/1.0.0/100"}), read only if network access is
  on.
\item[onWarning] A function that receives the warnings (without the \code{WFDB-x: }
  prefix and the newline). Default \code{null}: the library prints \code{WFDB-x: }
  followed by the warning on the standard error, unless \code{wfdbquiet()} was called.
\item[inputDir] The input directory for this call only (optional): a local directory, a
  \code{\textasciitilde} path or a URL, under which relative record names are read. It
  does not apply to full paths or URLs. Default \code{null}: the current setting
  (\code{wfdbxSetInputDir}).
\item[network] \code{true} to allow reading URLs for this call only, \code{false} to
  forbid it. Default \code{null}: the current setting.
\end{params}
\returns{A list of (key, value) pairs, in file order. A key that appears more than once
(ignoring letter case) is returned once, with its first value, and a warning is issued
(\cref{sec:kotlin-warnings}).}
\raises{\code{WfdbException} if the header cannot be found or read (``can't read header
for record REC''); \code{WfdbArgumentException} if \code{record} is empty or \code{"-"}
(the standard input).}
\begin{lstlisting}[language=Kotlin,style=example]
import wfdbx.*

writeRecord("pat", IntArray(360), nsig = 1, fs = 360.0, gains = doubleArrayOf(200.0),
            adcres = intArrayOf(12), units = listOf("mV"), names = listOf("ECG"),
            comments = listOf("recorded at rest"), metadata = listOf("Age" to "63", "Sex" to "F"))
println(readMetadata("pat"))           // [(Age, 63), (Sex, F)]
println(readHeader("pat").comments)
// [recorded at rest,  Age: 63,  Sex: F]: the free comment is not an entry
\end{lstlisting}

\apientry{getMetadata}{Return the value of a metadata entry}
\begin{lstlisting}[language=Kotlin,style=signature]
fun Wfdb.getMetadata(
    record: String,
    key: String,
    onWarning: ((String) -> Unit)? = null,
    inputDir: String? = null,
    network: Boolean? = null,
): String?
\end{lstlisting}
Returns the value of one metadata entry of the header of a record. The key is checked
before the header is read.
\begin{params}
\item[record] Record name (as for \code{readMetadata}).
\item[key] The key, such as \code{"Age"}. Letter case is ignored, and spaces at either end
  are removed.
\item[onWarning] A function that receives the warnings (as for \code{readMetadata}).
  Default \code{null}: the standard error.
\item[inputDir] The input directory for this call only. Default \code{null}: the current
  setting.
\item[network] \code{true} to allow reading URLs for this call only, \code{false} to
  forbid it. Default \code{null}: the current setting.
\end{params}
\returns{The value (\code{""} for an entry with an empty value), or \code{null} if the
header has no entry for \code{key}. If it has more than one, the first is returned, with a
warning.}
\raises{\code{WfdbException} if the header cannot be found or read;
\code{WfdbArgumentException} if \code{record} is empty or \code{"-"}, or if \code{key} is
not a valid key (see \code{setMetadata}).}
\begin{lstlisting}[language=Kotlin,style=example]
import wfdbx.*

writeRecord("pat", IntArray(360), nsig = 1, fs = 360.0, gains = doubleArrayOf(200.0),
            adcres = intArrayOf(12), units = listOf("mV"), metadata = listOf("Age" to "63", "Sex" to "F"))
println(getMetadata("pat", "age"))       // 63
println(getMetadata("pat", "Height"))    // null
\end{lstlisting}

\apientry{setMetadata}{Set the value of a metadata entry}
\begin{lstlisting}[language=Kotlin,style=signature]
fun Wfdb.setMetadata(
    record: String,
    key: String,
    value: String,
    inputDir: String? = null,
    network: Boolean? = null,
    dir: String? = null,
)
\end{lstlisting}
Sets the value of a metadata entry in the header of a record. If the header has an entry
for \code{key} (ignoring letter case), its line is replaced by \code{\# Key: value}, at
the same position and with the same line ending, with the spelling of \code{key} given
here (even if the value is the same); otherwise the entry is added after the last info
string of the header (or, if it has none, after its last line that is not blank), ending
as the record line ends. The entry is written as \code{\# Key: value}, or \code{\# Key:}
if the value is empty. Every other byte of the header is kept.
\begin{params}
\item[record] Record name (as for \code{readMetadata}). The header is changed where it was
  read, unless \code{dir} is given; a header read from a URL can only be written
  elsewhere, with \code{dir}.
\item[key] The key, such as \code{"Age"} or \code{"Date of birth"}. Spaces at either end
  are removed. It must not be empty, and must not contain \code{:}, \code{\#}, a line
  break or a NUL character.
\item[value] The value, which may be a sentence. Spaces at either end are removed. It
  must not contain a line break or a NUL character; it may be empty.
\item[inputDir] The input directory in which the header is read, for this call only.
  Default \code{null}: the current setting.
\item[network] \code{true} to allow reading the header from a URL, for this call only,
  \code{false} to forbid it. Default \code{null}: the current setting.
\item[dir] A directory, which must exist, in which to write the changed header, under the
  same file name, instead of changing the header itself (the header is then written even
  if nothing changes); the signal files are not copied. Default \code{null}: the header
  itself is changed.
\end{params}
\returns{Nothing (\code{Unit}).}
\raises{\code{WfdbArgumentException} if \code{key} or \code{value} is not valid
(``metadata key 'a:b' must not contain ':'{}'', ``value of metadata key 'Age' must not
contain a line break'', \ldots), if the header has more than one entry for \code{key}
(``header of record REC has more than one entry for key KEY (lines L1, L2); remove them
first''), if it has no record line, if it was read from a URL and \code{dir} is not given,
or if \code{record} is empty or \code{"-"}; \code{WfdbException} if the header cannot be
read (``can't read header for record REC''), if \code{dir} does not exist (``directory
DIR does not exist''), or if the changed header cannot be written (``can't write header
for record REC'').}
\begin{lstlisting}[language=Kotlin,style=example]
import wfdbx.*

writeRecord("pat", IntArray(360), nsig = 1, fs = 360.0, gains = doubleArrayOf(200.0),
            adcres = intArrayOf(12), units = listOf("mV"), metadata = listOf("Age" to "63", "Sex" to "F"))
setMetadata("pat", "AGE", "64")                                // replaces "# Age: 63"
setMetadata("pat", "Diagnosis", "sinus rhythm, no ectopy")     // a new entry
println(readMetadata("pat"))
// [(AGE, 64), (Sex, F), (Diagnosis, sinus rhythm, no ectopy)]

// Write the changed header to another directory; pat.hea is left unchanged.
val changed = wfdbxGetOutputDir() + "/changed"
java.io.File(changed).mkdir()                                  // (JVM)
setMetadata("pat", "Sex", "M", dir = changed)
println("${getMetadata("pat", "Sex")} ${getMetadata("$changed/pat", "Sex")}")    // F M
\end{lstlisting}

\apientry{removeMetadata}{Remove a metadata entry}
\begin{lstlisting}[language=Kotlin,style=signature]
fun Wfdb.removeMetadata(
    record: String,
    key: String,
    inputDir: String? = null,
    network: Boolean? = null,
    dir: String? = null,
): Int
\end{lstlisting}
Removes every entry for \code{key} (ignoring letter case) from the header of a record,
with its line ending, so it also repairs a header that has more than one. Other lines are
kept, byte for byte. Nothing is written if the header has no entry for \code{key}, unless
\code{dir} is given.
\begin{params}
\item[record] Record name (as for \code{setMetadata}).
\item[key] The key (as for \code{setMetadata}).
\item[inputDir] The input directory in which the header is read, for this call only.
  Default \code{null}: the current setting.
\item[network] \code{true} to allow reading the header from a URL, for this call only,
  \code{false} to forbid it. Default \code{null}: the current setting.
\item[dir] A directory, which must exist, in which to write the changed header instead of
  changing the header itself. Default \code{null}: the header itself is changed.
\end{params}
\returns{The number of lines removed (0 if the header has no entry for \code{key}).}
\raises{As \code{setMetadata}, except that more than one entry for \code{key} is not an
error.}
\begin{lstlisting}[language=Kotlin,style=example]
import wfdbx.*

writeRecord("pat", IntArray(360), nsig = 1, fs = 360.0, gains = doubleArrayOf(200.0),
            adcres = intArrayOf(12), units = listOf("mV"), metadata = listOf("Age" to "63", "Sex" to "F"))
println(removeMetadata("pat", "sex"))     // 1
println(removeMetadata("pat", "Sex"))     // 0: there is no such entry now
println(readMetadata("pat"))              // [(Age, 63)]
\end{lstlisting}

\apientry{updateMetadata}{Make several changes to the metadata}
\begin{lstlisting}[language=Kotlin,style=signature]
fun Wfdb.updateMetadata(
    record: String,
    changes: List<Pair<String, String?>>,
    inputDir: String? = null,
    network: Boolean? = null,
    dir: String? = null,
)

fun Wfdb.updateMetadata(
    record: String,
    changes: Map<String, String?>,
    inputDir: String? = null,
    network: Boolean? = null,
    dir: String? = null,
)
\end{lstlisting}
Makes several changes to the metadata of a record, in one write of the header. Each
change is applied in order, as \code{setMetadata} does, or, for a \code{null} value, as
\code{removeMetadata} does. All the changes are checked first: if one is not valid, none
is made. The second overload takes the changes as a map, applied in its iteration order
(the order of insertion for the maps of \code{mapOf} and \code{linkedMapOf}).
\begin{params}
\item[record] Record name (as for \code{setMetadata}).
\item[changes] The changes: a list of (key, value) pairs, or a map from key to value; a
  \code{null} value removes the entry. No two keys may differ only in letter case.
\item[inputDir] The input directory in which the header is read, for this call only.
  Default \code{null}: the current setting.
\item[network] \code{true} to allow reading the header from a URL, for this call only,
  \code{false} to forbid it. Default \code{null}: the current setting.
\item[dir] A directory, which must exist, in which to write the changed header instead of
  changing the header itself. Default \code{null}: the header itself is changed.
\end{params}
\returns{Nothing (\code{Unit}).}
\raises{As \code{setMetadata}; also \code{WfdbArgumentException} if two keys of
\code{changes} differ only in letter case (``metadata key 'dx' is given more than once
(keys are compared without regard to letter case)'').}
\begin{lstlisting}[language=Kotlin,style=example]
import wfdbx.*

writeRecord("pat", IntArray(360), nsig = 1, fs = 360.0, gains = doubleArrayOf(200.0),
            adcres = intArrayOf(12), units = listOf("mV"), metadata = listOf("Age" to "63", "Sex" to "F"))
updateMetadata("pat", mapOf("Age" to "64", "Sex" to null, "Dx" to "164889003"))
println(readMetadata("pat"))         // [(Age, 64), (Dx, 164889003)]
try {
    updateMetadata("pat", listOf("Dx" to "426783006", "dx" to null))
} catch (e: WfdbArgumentException) {
    println(e.message)        // metadata key 'dx' is given more than once (...)
}
println(getMetadata("pat", "Dx"))    // 164889003: nothing was changed
\end{lstlisting}

% ------------------------------------------------------------------------------------------
\subsection{Settings}
\label{sec:kotlin-settings}

These functions set where files are read and written, network access and legacy mode, on
a \code{Wfdb} object (the top-level functions: on the default instance), for the
high-level functions and the library functions alike (\cref{chap:files,chap:behavior}).
Each \code{set} function returns the previous setting, so that it can be restored. They
are the high-level forms of the library functions \code{wfdbx\_set\_input\_dir},
\code{wfdbx\_set\_network}, \ldots{} (\cref{sec:kotlin-library}), which take and return
bytes and integers.

\apientry{wfdbxSetInputDir}{Set the input directory}
\begin{lstlisting}[language=Kotlin,style=signature]
fun Wfdb.wfdbxSetInputDir(dir: String?): String?
\end{lstlisting}
Sets the input directory, from which names that are not full paths are read.
\begin{params}
\item[dir] A local directory, a \code{\textasciitilde} path or a URL; \code{null} (or
  \code{""}) restores the default: \code{WFDBX\_INPUT\_DIR}, or else the current directory.
\end{params}
\returns{The previous input directory, or \code{null} if none was set.}
\raises{Nothing.}
\begin{lstlisting}[language=Kotlin]
val previous = wfdbxSetInputDir("/data/mitdb")
val rec = readRecord("100")             // reads /data/mitdb/100.hea and 100.dat
wfdbxSetInputDir(previous)
\end{lstlisting}

\apientry{wfdbxGetInputDir}{Return the input directory}
\begin{lstlisting}[language=Kotlin,style=signature]
fun Wfdb.wfdbxGetInputDir(): String?
\end{lstlisting}
Returns the input directory. The function has no parameters.
\returns{The input directory, or \code{null} if none is set (the default).}
\raises{Nothing.}
\begin{lstlisting}[language=Kotlin]
println(Wfdb().wfdbxGetInputDir())      // null, unless WFDBX_INPUT_DIR is set
\end{lstlisting}

\apientry{wfdbxSetOutputDir}{Set the output directory}
\begin{lstlisting}[language=Kotlin,style=signature]
fun Wfdb.wfdbxSetOutputDir(dir: String?): String?
\end{lstlisting}
Sets the output directory, in which files named without a full path are created.
\begin{params}
\item[dir] A local directory or a \code{\textasciitilde} path (a URL is refused, with a
  message); \code{null} (or \code{""}) restores the default: \code{WFDBX\_OUTPUT\_DIR}, or
  else the current directory.
\end{params}
\returns{The previous output directory, or \code{null} if none was set.}
\raises{Nothing.}
\begin{lstlisting}[language=Kotlin]
wfdbxSetOutputDir("out")
writeAnnotations("100", "qrs", longArrayOf(100, 460), symbols = listOf("N", "N"))   // out/100.qrs
\end{lstlisting}

\apientry{wfdbxGetOutputDir}{Return the output directory}
\begin{lstlisting}[language=Kotlin,style=signature]
fun Wfdb.wfdbxGetOutputDir(): String?
\end{lstlisting}
Returns the output directory. The function has no parameters.
\returns{The output directory, or \code{null} if none is set (the default).}
\raises{Nothing.}
\begin{lstlisting}[language=Kotlin]
val out = wfdbxGetOutputDir() ?: "."
\end{lstlisting}

\apientry{wfdbxSetNetwork}{Turn network access on or off}
\begin{lstlisting}[language=Kotlin,style=signature]
fun Wfdb.wfdbxSetNetwork(on: Boolean): Boolean
\end{lstlisting}
Turns network access (reading URLs) on or off. It is off unless the environment variable
\code{WFDBX\_NETWORK} is a nonzero number.
\begin{params}
\item[on] \code{true} to allow reading URLs, \code{false} to forbid it.
\end{params}
\returns{The previous setting.}
\raises{Nothing.}
\begin{lstlisting}[language=Kotlin]
wfdbxSetNetwork(true)
wfdbxSetInputDir("https://physionet.org/files/mitdb/1.0.0/")
val rec = readRecord("100", to = TimeSpec("0:05"))      // from PhysioNet
\end{lstlisting}

\apientry{wfdbxGetNetwork}{Return whether network access is on}
\begin{lstlisting}[language=Kotlin,style=signature]
fun Wfdb.wfdbxGetNetwork(): Boolean
\end{lstlisting}
Returns whether network access is on. The function has no parameters.
\returns{\code{true} if network access is on.}
\raises{Nothing.}
\begin{lstlisting}[language=Kotlin]
if (!wfdbxGetNetwork()) println("network access is off")
\end{lstlisting}

\apientry{wfdbxSetLegacy}{Turn legacy mode on or off}
\begin{lstlisting}[language=Kotlin,style=signature]
fun Wfdb.wfdbxSetLegacy(on: Boolean): Boolean
\end{lstlisting}
Turns legacy mode (the behavior of the original WFDB package, including its assumed
signal parameters and its search of the WFDB path, \cref{chap:behavior}) on or off, for
the whole \code{Wfdb} object. The parameter \code{legacy} of the high-level functions
turns it on for one call.
\begin{params}
\item[on] \code{true} to turn legacy mode on, \code{false} to turn it off.
\end{params}
\returns{The previous setting.}
\raises{Nothing.}
\begin{lstlisting}[language=Kotlin]
val previous = wfdbxSetLegacy(true)
val rec = readRecord("100s")            // as the original package reads it
wfdbxSetLegacy(previous)
\end{lstlisting}

\apientry{wfdbxGetLegacy}{Return whether legacy mode is on}
\begin{lstlisting}[language=Kotlin,style=signature]
fun Wfdb.wfdbxGetLegacy(): Boolean
\end{lstlisting}
Returns whether legacy mode is on. The function has no parameters.
\returns{\code{true} if legacy mode is on.}
\raises{Nothing.}
\begin{lstlisting}[language=Kotlin]
println(wfdbxGetLegacy())               // false, unless it was turned on
\end{lstlisting}

% ------------------------------------------------------------------------------------------
\subsection{The result classes}
\label{sec:kotlin-results}

The readers return objects of the classes \code{Record}, \code{Header} and
\code{Annotations}, whose properties are listed in
\cref{tab:kotlin-record,tab:kotlin-header,tab:kotlin-annotations}; \code{checkRecord}
returns a \code{CheckReport}, which holds \code{Issue} objects
(\cref{tab:kotlin-checkreport,tab:kotlin-issue}). The signals of a record are described
by \code{Siginfo} objects (\cref{tab:kotlin-siginfo}), which are those of the WFDB library
(\cref{sec:kotlin-library}). \code{Record}, \code{Header}, \code{Annotations},
\code{CheckReport}, \code{Issue} and \code{BaseDate} are data classes, whose properties
are read-only (\code{val}); \code{Record} and \code{Annotations} compare their arrays by
content in \code{equals}. A \code{BaseDate} has the properties \code{year}, \code{month}
(1--12) and \code{day} (1--31).

\begin{xltabular}{\linewidth}{@{}lL@{}}
\caption{The properties and methods of \code{Record}, the result of \code{readRecord}}
\label{tab:kotlin-record} \\
\toprule
\tabhead Property & \tabhead Contents \\
\midrule
\endfirsthead
\toprule
\tabhead Property & \tabhead Contents \\
\midrule
\endhead
\bottomrule
\endfoot
\code{name} & the record name, as given to \code{readRecord} (\code{String}) \\
\code{fs} & the sampling frequency in Hz (\code{Double}), or 0 if the header gives none;
  with \code{highres = true}, the highest sampling frequency of a multi-frequency
  record \\
\code{nsig} & the number of signals returned \\
\code{nsamples} & the number of samples per signal (the number of frames, or of sample
  vectors with \code{highres}): \code{samples.size / nsig} \\
\code{signalNumbers} & the number, in the record, of each signal returned
  (\code{IntArray}) \\
\code{samples} & the samples in ADC units, frame by frame (\code{IntArray}):
  \code{samples[i * nsig + s]} is sample \code{i} of signal \code{s}; invalid samples are
  \code{WFDB\_INVALID\_SAMPLE} \\
\code{physical} & with \code{physical = true}, the samples in physical units, in the
  layout of \code{samples} (\code{DoubleArray}; NaN for an invalid sample); otherwise
  \code{null} \\
\code{siginfo} & a \code{List<Siginfo>}, one per signal returned; \code{desc} has no white
  space at either end, and \code{adcres} is 0 if the header does not give it \\
\code{names} & the description of each signal, such as \code{"MLII"} (\code{""} if none),
  without the white space that a header may put at either end of it \\
\code{units} & the physical units of each signal, or \code{null} where the header gives
  none (\code{List<String?>}) \\
\code{gains} & the gain of each signal, in ADC units per physical unit
  (\code{List<Double>}); 0 means that the signal is uncalibrated \\
\code{baselines} & the ADC value of each signal that corresponds to 0 physical units
  (\code{List<Int>}) \\
\code{baseTime} & the start time and date as text, to the second, such as
  \code{"[14:00:03 21/06/2024]"}, or \code{null} if the header gives none \\
\code{baseTimeMs} & the start time in milliseconds after midnight, such as
  \code{50403500} for 14:00:03.500 (\code{Long}), or \code{null} if \code{baseTime} is
  \code{null} \\
\code{baseDate} & the start date as a \code{BaseDate(year, month, day)}, or \code{null}
  if the header gives none \\
\code{sampfrom} & the number of the first sample returned (\code{Long}) \\
\code{comments} & the info strings of the header (without the \code{\#}), and of the
  record's \file{.info} file, if any \\
\code{checksumOk} & \code{false} if the samples do not match the checksums of the
  header; checksums are checked only when a read reaches the end of the record, so a
  shorter read gives \code{true} \\
\code{legacy} & \code{true} if the record was read in legacy mode (a gain of 200 is then
  used for uncalibrated signals in physical units) \\
\code{at(i, s)} & sample \code{i} of signal \code{s}, in ADC units \\
\code{signal(s)} & all the samples of signal \code{s}, in ADC units (\code{IntArray});
  \code{IllegalArgumentException} (``no signal S'') if there is no signal \code{s} \\
\code{physicalAt(i, s)} & sample \code{i} of signal \code{s} in physical units,
  $(\mathit{adc} - \mathit{baseline}) / \mathit{gain}$ (NaN for an invalid sample);
  \code{WfdbException} (``signal N of record REC is uncalibrated'') if the signal has no
  gain \\
\code{physicalSignal(s)} & all the samples of signal \code{s} in physical units
  (\code{DoubleArray}), as \code{physicalAt} \\
\end{xltabular}

\begin{xltabular}{\linewidth}{@{}lL@{}}
\caption{The properties of \code{Header}, the result of \code{readHeader}}
\label{tab:kotlin-header} \\
\toprule
\tabhead Property & \tabhead Contents \\
\midrule
\endfirsthead
\toprule
\tabhead Property & \tabhead Contents \\
\midrule
\endhead
\bottomrule
\endfoot
\code{name} & the record name, as given to \code{readHeader} \\
\code{fs} & the sampling frequency in Hz, or 0 if the header gives none \\
\code{counterFreq} & the counter frequency, as \code{getcfreq()} returns it \\
\code{baseCounter} & the base counter value, as \code{getbasecount()} returns it \\
\code{nsamples} & the number of samples per signal (\code{Long}), or 0 if the header does
  not give it \\
\code{nsig} & the number of signals \\
\code{siginfo} & a \code{List<Siginfo>}, one per signal; \code{desc} has no white space at
  either end, and \code{adcres} is 0 if the header does not give it \\
\code{baseTime}, \code{baseTimeMs}, \code{baseDate} & the start time and date, as in
  \code{Record} \\
\code{comments} & the info strings of the header (without the \code{\#}) \\
\code{metadata} & the (key, value) pairs of the metadata entries of the header
  (\code{List<Pair<String, String>{}>}), in file order, as \code{readMetadata} returns
  them; empty for an EDF file or a header read from the standard input \\
\end{xltabular}

\begin{xltabular}{\linewidth}{@{}lL@{}}
\caption{The properties of \code{Annotations}, the result of \code{readAnnotations}.
  The annotations are in file order; each array or list has one entry per annotation.}
\label{tab:kotlin-annotations} \\
\toprule
\tabhead Property & \tabhead Contents \\
\midrule
\endfirsthead
\toprule
\tabhead Property & \tabhead Contents \\
\midrule
\endhead
\bottomrule
\endfoot
\code{record} & the record name, as given \\
\code{annotator} & the annotator name, such as \code{"atr"} \\
\code{fs} & the frequency in which \code{samples} are counted, normally the sampling
  frequency of the record; 0 if neither the record, nor the file, nor the argument
  \code{fs} of \code{readAnnotations} gives one \\
\code{samples} & the annotation times, in samples (\code{LongArray}) \\
\code{codes} & the annotation codes (\code{IntArray}), such as \code{NORMAL} or
  \code{PVC} \\
\code{symbols} & the mnemonics of the codes, such as \code{"N"}, \code{"V"} or \code{"+"}
  (\code{List<String>}) \\
\code{subtypes} & the subtype of each annotation (\code{IntArray}) \\
\code{chans} & the channel number of each annotation (\code{IntArray}) \\
\code{nums} & the annotator number of each annotation (\code{IntArray}) \\
\code{aux} & the auxiliary text of each annotation, or \code{null} where there is none
  (\code{List<String?>}) \\
\code{size} & the number of annotations \\
\end{xltabular}

\begin{table}[htbp]
\centering
\caption{The properties and methods of \code{CheckReport}, the result of
  \code{checkRecord}}
\label{tab:kotlin-checkreport}
\begin{tabularx}{\linewidth}{@{}lL@{}}
\toprule
\tabhead Property & \tabhead Contents \\
\midrule
\code{record} & the record name, as given \\
\code{ok} & \code{true} if no issue has the severity \code{Severity.ERROR} \\
\code{issues} & the \code{List<Issue>}, in the order of the checks \\
\code{lines()} & the text form of the report, as a list of lines: one line per issue
  (\code{issue.toString()}), or the one line \code{no problems found} \\
\code{toString()} & the lines of \code{lines()}, joined by line feeds (what
  \code{println(report)} prints) \\
\bottomrule
\end{tabularx}
\end{table}

\begin{table}[htbp]
\centering
\caption{The properties of \code{Issue}, a problem found by \code{checkRecord}}
\label{tab:kotlin-issue}
\begin{tabularx}{\linewidth}{@{}lL@{}}
\toprule
\tabhead Property & \tabhead Contents \\
\midrule
\code{severity} & a \code{Severity}: \code{ERROR}, \code{WARNING} or \code{INFO}; its
  \code{text} (and \code{toString()}) is \code{"error"}, \code{"warning"} or
  \code{"info"} \\
\code{code} & a stable identifier of the kind of problem, such as \code{"checksum"} or
  \code{"gain\_missing"} (\cref{tab:python-issue-codes}) \\
\code{where} & where the problem is: \code{""} (the record), \code{"signal N"},
  \code{"segment SEG"}, \code{"segment SEG, signal N"} or \code{"annotator ANN"} \\
\code{message} & the description of the problem \\
\code{toString()} & \code{"SEVERITY: WHERE: MESSAGE"}, or \code{"SEVERITY: MESSAGE"} when
  \code{where} is empty \\
\bottomrule
\end{tabularx}
\end{table}
Two \code{Issue} objects are equal (\code{==}) when their four properties are. The codes
of the issues are those of \cref{tab:python-issue-codes}.

\begin{xltabular}{\linewidth}{@{}lL@{}}
\caption{The properties of \code{Siginfo}, the specification of a signal (mutable, with
  C's zero values by default)}
\label{tab:kotlin-siginfo} \\
\toprule
\tabhead Property & \tabhead Contents \\
\midrule
\endfirsthead
\toprule
\tabhead Property & \tabhead Contents \\
\midrule
\endhead
\bottomrule
\endfoot
\code{fname} & the name of the signal file, as the header gives it (\code{ByteArray?}) \\
\code{desc} & the signal description, such as \code{"MLII"} (\code{ByteArray?}) \\
\code{units} & the physical units, such as \code{"mV"}, or \code{null} if not specified
  (\code{ByteArray?}) \\
\code{gain} & the gain, in ADC units per physical unit (\code{Double}; 0:
  uncalibrated) \\
\code{initval} & the value of the first sample \\
\code{group} & the signal group number (the signals of one group share one file) \\
\code{fmt} & the storage format, such as \code{16}, \code{212} or \code{516} \\
\code{spf} & the number of samples per frame (more than 1 for the faster signals of a
  multi-frequency record) \\
\code{bsize} & the block size, in bytes (0 for ordinary files) \\
\code{adcres} & the ADC resolution, in bits (0: not specified) \\
\code{adczero} & the ADC output for an input of 0\,V \\
\code{baseline} & the ADC value that corresponds to 0 physical units \\
\code{nsamp} & the number of samples (\code{Long}; 0 if unknown) \\
\code{cksum} & the 16-bit checksum of the samples \\
\end{xltabular}

% ------------------------------------------------------------------------------------------
\subsection{Warnings and errors}
\label{sec:kotlin-warnings}

\paragraph{Exceptions.} The high-level functions report errors by exceptions
(\cref{tab:kotlin-exceptions}), whose messages say what failed. Each function checks its
arguments before it creates or changes a file, so an invalid argument leaves the files as
they were. The library may also print a message on the standard error when it cannot open
a file (as in \code{init: can't open header for record REC}); \code{wfdbquiet()} turns
these messages off (\cref{sec:kotlin-library}).

\begin{xltabular}{\linewidth}{@{}lL@{}}
\caption{The exceptions of the high-level API}
\label{tab:kotlin-exceptions} \\
\toprule
\tabhead Exception & \tabhead Thrown when \\
\midrule
\endfirsthead
\toprule
\tabhead Exception & \tabhead Thrown when \\
\midrule
\endhead
\bottomrule
\endfoot
\code{WfdbException} & every error of a high-level function: a header, signal file or
  annotation file cannot be found, read, created or written completely (for example on a
  full disk); a signal or annotation file ends before its header or its end-of-file marker
  says; an argument is invalid, or the operation needs a value that the record does not
  give (a time in seconds without a sampling frequency, physical units without a gain);
  the directory given as \code{dir} does not exist \\
\code{WfdbArgumentException} & a subclass of \code{WfdbException}, thrown by the metadata
  functions (and for the \code{metadata} argument of \code{writeRecord}) for an invalid
  argument: an invalid key or value, a key given twice, a header with two entries for the
  key to set, a header without a record line, a header read from a URL without
  \code{dir}, an empty record name or \code{"-"} (Python's \code{ValueError}); a header
  that cannot be read or written, or a \code{dir} that does not exist, is a plain
  \code{WfdbException} (Python's \code{IOError}) \\
\end{xltabular}

\paragraph{The checksum warning.} When \code{readRecord} reads to the end of a record
whose samples do not match the checksums of its header, it returns all the samples and
warns: ``checksum error in record REC: the signal data do not match the checksums in its
header''. \code{checksumWarning = false} turns the warning off; \code{checksumOk} of the
result is \code{false} either way.

\paragraph{The duplicate-key warning.} When a header holds two metadata entries whose keys
differ only in letter case (written by other software), \code{readMetadata},
\code{getMetadata} and \code{readHeader} use the first and warn: ``header of record REC
has more than one entry for key KEY (lines L1, L2); the first is used''.
\code{setMetadata} and \code{updateMetadata} then refuse to change that key
(\code{WfdbArgumentException}), and \code{removeMetadata} removes all of its entries.

Both warnings go to the function given as \code{onWarning}, which receives the message;
by default, the library prints \code{WFDB-x: } followed by the message on the standard
error, through \code{wfdb\_error}, so that \code{wfdbquiet()} silences it and
\code{wfdberror()} returns it:

\begin{lstlisting}[language=Kotlin,style=example]
import wfdbx.*

// A record whose last sample is changed after it was written: the checksum fails.
writeRecord("ck", IntArray(100) { it }, nsig = 1, fs = 100.0, gains = doubleArrayOf(100.0),
            adcres = intArrayOf(12), units = listOf("mV"))
java.io.RandomAccessFile(wfdbxGetOutputDir() + "/ck.dat", "rw").use { f ->
    f.seek(f.length() - 2)
    f.write(byteArrayOf(0, 0))
}
val caught = ArrayList<String>()
var rec = readRecord("ck", onWarning = { caught.add(it) })
println("${rec.nsamples} ${rec.checksumOk}")    // 100 false: all the samples
println(caught[0])
// checksum error in record ck: the signal data do not match the checksums in its header
rec = readRecord("ck", checksumWarning = false)      // no warning

// Two entries for one key: the first is used, with a warning.
writeRecord("dup", IntArray(10), nsig = 1, fs = 100.0, gains = doubleArrayOf(0.0),
            adcres = intArrayOf(12), comments = listOf(" Age: 63", " AGE: 64"))
println(getMetadata("dup", "age", onWarning = { caught.add(it) }))     // 63
println(caught[1])
// header of record dup has more than one entry for key Age (lines 3, 4); the first is used
try {
    setMetadata("dup", "Age", "65")
} catch (e: WfdbArgumentException) {
    println("WfdbArgumentException: ${e.message}")
}
println(removeMetadata("dup", "age"))       // 2
\end{lstlisting}

\begin{lstlisting}[language=Kotlin,style=example]
import wfdbx.*

wfdbquiet()                           // no library messages on the standard error
try {
    readRecord("nosuchrecord")
} catch (e: WfdbException) {
    println("WfdbException: ${e.message}")    // WfdbException: can't read record nosuchrecord
}
wfdbverbose()

writeRecord("raw", IntArray(10) { it }, nsig = 1, fs = 100.0, gains = doubleArrayOf(0.0),
            adcres = intArrayOf(12))          // uncalibrated
try {
    readRecord("raw", physical = true)
} catch (e: WfdbException) {
    println("WfdbException: ${e.message}")    // WfdbException: signal 0 of record raw is uncalibrated
}
\end{lstlisting}

% ==========================================================================================
\section{The WFDB library functions in Kotlin}
\label{sec:kotlin-library}

Every function of the WFDB library is available under its C name, with the same
arguments, in the same order, and the same return values, so that programs and
documentation written for the C library carry over directly. The functions, their
arguments and their results are documented once, for all the languages, in
\cref{chap:library}; this section gives the rules by which the C declarations become
Kotlin functions. The KDoc comment of each function documents it, including what WFDB-x
changes in it.

\subsection{Conventions}

\paragraph{The library state.} All the state that the C library keeps in global variables
(the WFDB path, the open record, its signals, the annotators, the calibration list, the
settings) is held by a \code{Wfdb} object. Each C function is an extension function of
\code{Wfdb} with the C name and arguments, \code{fun Wfdb.isigopen(record, siarray,
nsig)}, called as \code{w.isigopen(...)}, and also a top-level function with the same name
and arguments, \code{isigopen(...)}, which uses the default instance of the calling thread
(\code{wfdbDefault()}), for line-by-line ports of C programs. Separate \code{Wfdb} objects
share no state, so several records can be open at once, in one thread or in several; a
\code{Wfdb} object is used by one thread at a time (it may be handed from one thread to
another). As in C, the open record, its input signals, the output signals and the
annotators remain open between calls, until \code{wfdbquit()} closes them;
\code{wfdbquit()} leaves the input and output directories and the network setting as they
are. There are no destructors: an output file is completed by \code{wfdbquit()} (or the
function that closes it), and \code{wfdbx.cstd.exit} flushes the open output files as C's
\code{exit} does, but writes no annotation end marks and no header updates.

\paragraph{Names.} The functions, the constants (\code{WFDB\_READ},
\code{WFDB\_INVALID\_SAMPLE}, \code{WFDB\_HIGHRES}, \ldots) and the annotation codes
(\code{NORMAL}, \code{PVC}, \ldots, \cref{chap:formats}) keep their C names. The functions
of \file{<wfdb/ecgmap.h>} are provided under their macro names (\code{isqrs},
\code{map1}, \code{map2}, \ldots).

\paragraph{Strings are bytes.} WFDB files are not UTF-8: every string of the C-style API
is a \code{ByteArray}, which, like a C string, ends at its first NUL byte, and C's
\code{NULL} is \code{null}. Where natural, a function also has an overload that takes a
\code{String}, one character per byte (ISO-8859-1): \code{"..".latin1()} and
\code{bytes.latin1()}, from \code{wfdbx.cstd}, convert. Where both overloads exist, a
\code{null} argument is written \code{null as ByteArray?}, as in
\code{sampfreq(null as ByteArray?)}. Functions that return a C string (\code{timstr},
\code{mstimstr}, \code{datstr}, \code{getinfo}, \ldots) return a new \code{ByteArray} (a
copy of C's static buffer); \code{getinfo} returns \code{null} after the last info
string.

\paragraph{Numbers.} C's \code{int} and \code{WFDB\_Sample} are \code{Int},
\code{long}, \code{WFDB\_Time} and \code{WFDB\_Date} are \code{Long}, and \code{double}
and \code{WFDB\_Frequency} are \code{Double}. \code{WFDB\_Signal} and \code{WFDB\_Group}
are \code{Int}; negative values are invalid signal numbers, as C's unsigned values are.
The \code{char} fields of the structures are \code{Int} properties that keep C's range:
\code{anntyp}, \code{subtyp} and \code{num} from $-128$ to 127, \code{chan} from 0 to
255.

\paragraph{Arrays and output arguments.} Where a C function fills an array that the caller
passes, the Kotlin function fills an array that the caller passes:
\begin{itemize}
\item \code{isigopen(record, siarray, nsig)} (also \code{osigopen} and \code{wfdbinit})
  takes an \code{Array<Siginfo>?}, whose elements it fills (allocated by the caller:
  \code{Array(n) \{ Siginfo() \}}); \code{null} is C's \code{NULL}. With \code{nsig}
  $\le 0$, \code{isigopen} reads the header only, and returns the number of signals of the
  record.
\item \code{getvec(v)}, \code{getframe(v)} and \code{putvec(v)} take an \code{IntArray};
  a shorter array receives what fits.
\item A C function that fills a structure through a pointer fills the object passed:
  \code{getann(n, annot)} fills an \code{Annotation}, \code{getcal(desc, units, cal)} a
  \code{Calinfo}. \code{putann(n, annot)} and \code{putcal(cal)} read them.
\item An array of structures passed in is an array of objects, as in
  \code{osigfopen(siarray, nsig)} and \code{annopen(record, aiarray, nann)}, where
  \code{aiarray} is, for example, \code{arrayOf(Anninfo("atr", WFDB\_READ))}.
\item \code{getseginfo(sarray)} takes a \code{MutableList<Seginfo>?}, which it clears and
  fills with copies of the segments; \code{setmsheader(record, names, n)} takes an
  \code{Array<ByteArray>} or an \code{Array<String>}.
\end{itemize}

\paragraph{Return values and errors.} The functions return what the C functions return:
an \code{Int} status (negative on errors), a number of signals, a \code{Double}, a byte
string or \code{null}. Errors are reported by these values and by a message on the
standard error, as in C, never by exceptions. \code{wfdberror()} returns the last message;
\code{wfdbquiet()} turns the messages off and \code{wfdbverbose()} on again.
\code{wfdbx\_write\_errors()} returns the number of write errors since the program
started (output files that could not be written, flushed or closed),
\code{wfdbx\_note\_write\_error()} counts one more, and
\code{wfdbx\_fclose\_output(file, pname, name)} closes a file that a program writes itself,
reporting ``PNAME: write error on NAME'' and counting it.
\code{wfdbx\_input\_path(name)} and \code{wfdbx\_output\_path(name)} return the path under
which a program opens a file that it reads or creates itself, where the library would
open a file of that name: under the input or output directory (the missing directories
of an output path are created), a \code{\textasciitilde} name expanded, a full path as
given; \code{null} for a URL, and \code{"-"} as given. In legacy mode, an input name is
used as given, and the output directory, if set, is a prefix of any output name.
\code{wfdbquit()} closes the open record and annotators, and keeps the directory and
network settings, the time resolution of output annotators (\code{getafreq()}) and the
getvec mode (\code{getgvmode()}), as in C; the high-level functions restore the last two
(\cref{sec:kotlin-conventions}).

\paragraph{Additions.} \code{getvecs(out, n)} reads up to \code{n} sample vectors at once
into the \code{IntArray} \code{out}, vector after vector, as \code{n} calls of
\code{getvec} would, but decoding whole blocks of the signal files; it returns the number
of vectors read, or the negative status of \code{getvec} if the first read failed.
\code{getvecs\_status(out, n)} returns a \code{GetvecsResult(count, status)}, with the
value returned by the last \code{getvec}. The C standard library functions that the WFDB
library relies on are in \code{wfdbx.cstd}: \code{cformat} and \code{sprintf} (to a
\code{ByteArray}), \code{printf} and \code{eprintf}, which format numbers exactly as the
GNU C library does, \code{strtol}, \code{strtod}, \code{sscanf}, \ldots

\paragraph{Settings.} The settings functions of WFDB-x (\code{wfdbx\_set\_input\_dir},
\code{wfdbx\_get\_input\_dir}, \code{wfdbx\_set\_output\_dir},
\code{wfdbx\_get\_output\_dir}, \code{wfdbx\_set\_network}, \code{wfdbx\_get\_network},
\code{wfdbx\_set\_legacy}, \code{wfdbx\_get\_legacy}, \code{wfdbx\_legacy\_args}) keep
their C names and types (byte strings, \code{Int} for the switches); the high-level
functions of \cref{sec:kotlin-settings} are their Kotlin forms.

\subsection{Package layout}

\begin{table}[htbp]
\centering
\caption{The packages of the Kotlin library}
\label{tab:kotlin-packages}
\begin{tabularx}{\linewidth}{@{}lL@{}}
\toprule
\tabhead Package or file & \tabhead Contents \\
\midrule
\code{wfdbx}: \file{Wfdb.kt} & the class \code{Wfdb}, \code{wfdbDefault()},
  \code{withDefault} \\
\file{Types.kt} & the constants and structures of \file{wfdb.h} \\
\file{WfdbFile.kt}, \file{Wfdbio.kt} & files, the input and output directories, network
  access, error messages (\file{lib/wfdbio.c}) \\
\file{Signal.kt}, \file{SignalTime.kt}, \file{Flac.kt} & headers, signal files in every
  format, time and date conversions (\file{lib/signal.c}), the FLAC codec \\
\file{Annot.kt}, \file{Calib.kt}, \file{Wfdbinit.kt}, \file{Ecgcodes.kt} & annotation files
  (\file{lib/annot.c}), the calibration list (\file{lib/calib.c}), \code{wfdbinit}, the
  annotation codes and \file{ecgmap.h} \\
\file{Record.kt}, \file{Check.kt}, \file{Metadata.kt} & the high-level API
  (\cref{sec:kotlin-api}) \\
\code{wfdbx.cstd} & the C standard library emulation (\code{printf}, \code{scanf},
  \code{strtol}, \code{rand}, the standard streams, \code{exit}, \code{latin1}) and the
  platform layer \\
\code{wfdbx.apps} & the command-line tools, one file per tool, as functions
  \code{fun Wfdb.rdsampMain(argv: Array<ByteArray>): Int}, \ldots \\
\code{wfdbx.cli} & the command-line driver \code{wfdbx.cli.MainKt} (JVM) \\
\bottomrule
\end{tabularx}
\end{table}

\subsection{The structures}
\label{sec:kotlin-structures}

The structures of the C library are mutable classes whose properties have the names of
the fields of the C structures and C's zero values by default
(\cref{tab:kotlin-structures}). \code{copy()} is C's structure assignment, and
\code{set(o)} copies another object into an existing one (what a C function filling
\code{*p} does); \code{equals} compares contents.
\begin{lstlisting}[language=Kotlin]
Siginfo(fname: ByteArray? = null, desc: ByteArray? = null, units: ByteArray? = null, gain: Double = 0.0,
        initval: Int = 0, group: Int = 0, fmt: Int = 0, spf: Int = 0, bsize: Int = 0, adcres: Int = 0,
        adczero: Int = 0, baseline: Int = 0, nsamp: Long = 0, cksum: Int = 0)
Seginfo(recname: ByteArray = ByteArray(0), nsamp: Long = 0, samp0: Long = 0)
Anninfo(name: ByteArray = ByteArray(0), stat: Int = 0)      // also Anninfo(name: String, stat: Int)
Annotation(time: Long = 0, anntyp: Int = 0, subtyp: Int = 0, chan: Int = 0, num: Int = 0, aux: ByteArray? = null)
Calinfo(low: Double = 0.0, high: Double = 0.0, scale: Double = 0.0, sigtype: ByteArray? = null,
        units: ByteArray? = null, caltype: Int = 0)
\end{lstlisting}

\begin{xltabular}{\linewidth}{@{}lL@{}}
\caption{The classes of the C structures}
\label{tab:kotlin-structures} \\
\toprule
\tabhead Class (C structure) & \tabhead Properties \\
\midrule
\endfirsthead
\toprule
\tabhead Class (C structure) & \tabhead Properties \\
\midrule
\endhead
\bottomrule
\endfoot
\code{Siginfo} (\code{WFDB\_Siginfo}) & \code{fname}, \code{desc}, \code{units},
  \code{gain}, \code{initval}, \code{group}, \code{fmt}, \code{spf}, \code{bsize},
  \code{adcres}, \code{adczero}, \code{baseline}, \code{nsamp}, \code{cksum}
  (\cref{tab:kotlin-siginfo}) \\
\code{Seginfo} (\code{WFDB\_Seginfo}) & \code{recname} (the name of the segment),
  \code{nsamp} (its number of samples), \code{samp0} (the number of its first sample in
  the record) \\
\code{Anninfo} (\code{WFDB\_Anninfo}) & \code{name} (the annotator name, the suffix of the
  file), \code{stat} (\code{WFDB\_READ} or \code{WFDB\_WRITE} for the MIT format,
  \code{WFDB\_AHA\_READ} or \code{WFDB\_AHA\_WRITE} for the AHA format) \\
\code{Annotation} (\code{WFDB\_Annotation}) & \code{time} (in sample intervals),
  \code{anntyp} (the annotation code), \code{subtyp}, \code{chan}, \code{num},
  \code{aux} (\code{null}, or bytes whose first byte is the length of the text that
  follows, as in C) \\
\code{Calinfo} (\code{WFDB\_Calinfo}) & \code{low}, \code{high} (the levels of the
  calibration pulse, in physical units), \code{scale} (physical units per cm),
  \code{sigtype} (matched against the beginning of a signal description), \code{units},
  \code{caltype} \\
\end{xltabular}

\subsection{Examples}

Reading the samples and the annotations of a record, as a C program does, on a
\code{Wfdb} object of its own (the input directory is that of the examples):

\begin{lstlisting}[language=Kotlin,style=example]
import wfdbx.*
import wfdbx.cstd.latin1

val w = Wfdb()                              // a library state of its own
w.wfdbx_set_input_dir(wfdbxGetInputDir()!!)
val si = Array(2) { Siginfo() }             // filled with the signal specifications
val nsig = w.isigopen("100s", si, 2)        // open 2 signals
println("$nsig ${si.map { "${it.desc?.latin1()} ${it.fmt} ${it.gain}" }}")   // 2 [MLII 212 200.0, V5 212 200.0]
println(w.sampfreq(null as ByteArray?))     // 360.0: the frequency of the open record

val v = IntArray(nsig)                      // receives one sample per signal
w.isigsettime(w.strtim("0:30"))             // skip to 30 seconds
w.getvec(v)
println("${v.toList()} ${(0 until nsig).map { w.aduphys(it, v[it]) }}")   // [947, 968] [-0.385, -0.28]

val frames = IntArray(360 * nsig)
println(w.getvecs(frames, 360))             // 360: the next second, frame by frame

val a = Annotation()
w.annopen("100s", arrayOf(Anninfo("atr", WFDB_READ)), 1)
var n = 0
while (w.getann(0, a) == 0) {               // 0: an annotation was read; -1: the end
    if (n < 2) println("${w.mstimstr(a.time).latin1().trim()} ${w.annstr(a.anntyp).latin1()}")   // 0:00.050 +, 0:00.214 N
    n++
}
println("$n annotations")                   // 75 annotations

w.wfdbquiet()                               // no messages on the standard error
println("${w.isigopen("nosuchrecord", null, 0)} ${w.wfdberror().latin1().trim()}")
// -1 init: can't open header for record nosuchrecord
w.wfdbverbose()
w.wfdbquit()                                // close everything
\end{lstlisting}

Writing a record and an annotation file, sample by sample and annotation by annotation:

\begin{lstlisting}[language=Kotlin,style=example]
import wfdbx.*
import wfdbx.cstd.latin1

val w = Wfdb()
w.wfdbx_set_output_dir(wfdbxGetOutputDir()!!)
val out = arrayOf(Siginfo(fname = "ramp.dat".latin1(), desc = "ramp".latin1(), units = "mV".latin1(),
                          gain = 200.0, fmt = 16, adcres = 12))
println(w.osigfopen(out, 1))                // 1: one output signal
w.setsampfreq(250.0)
for (k in 0 until 500) w.putvec(intArrayOf(k))
println(w.newheader("ramp"))                // 0: ramp.hea is written
w.annopen("ramp", arrayOf(Anninfo("qrs", WFDB_WRITE)), 1)
for (t in longArrayOf(100, 350)) w.putann(0, Annotation(time = t, anntyp = NORMAL))
w.wfdbquit()                                // completes ramp.dat and ramp.qrs

println(readAnnotations("ramp", "qrs").samples.toList())    // [100, 350]
\end{lstlisting}

% ==========================================================================================
\section{The command-line tools in Kotlin}
\label{sec:kotlin-tools}

The 33 command-line tools (\code{rdsamp}, \code{rdann}, \code{wrsamp}, \code{xform},
\ldots) run on the JVM, through the launcher \file{kotlin/bin/wfdbx} (a POSIX shell
script), followed by the name of the tool and its arguments. The tools, their options,
their output and their exit status are those of the C tools, and are documented in
\cref{chap:tools}. The launcher needs the jar built by \cmd{./gradlew fatJar}
(\cref{sec:kotlin-install}). In the directory \file{kotlin}:

\begin{lstlisting}[style=shell]
export WFDBX_INPUT_DIR=../tests/data
bin/wfdbx rdsamp -r 100s -t 10 -p      # the first 10 seconds of 100s, in physical units
bin/wfdbx rdann -r 100s -a atr         # its annotations
bin/wfdbx rdsamp --legacy -r 100s      # legacy mode: --legacy as the first argument
\end{lstlisting}

The command also takes these options, instead of a tool:

\begin{table}[htbp]
\centering
\caption{The options of the launcher \code{bin/wfdbx}}
\label{tab:kotlin-wfdbx-options}
\begin{tabularx}{\linewidth}{@{}lL@{}}
\toprule
\tabhead Command & \tabhead Effect \\
\midrule
\cmd{bin/wfdbx TOOL [ARGUMENTS ...]} & runs the tool \code{TOOL}, and exits with its
  status \\
\cmd{bin/wfdbx --list} & lists the names of the tools, one per line \\
\cmd{bin/wfdbx --version} & prints the version, \code{WFDB-x 1.2.0} \\
\cmd{bin/wfdbx --help}, \cmd{bin/wfdbx -h} & prints the usage \\
\bottomrule
\end{tabularx}
\end{table}

Without arguments, the command prints the usage and exits with status 1; an unknown tool
is an error (``wfdbx: unknown tool 'TOOL' (try 'wfdbx --list')'', status 1), and an
uncaught exception is reported as \code{wfdbx: error: MESSAGE} (status 1). When an output
file or the standard output cannot be written (a full disk, for example), a tool reports
it and exits with status 1 (legacy mode keeps the original status, 0).

The files that a tool reads or creates itself are found where the library finds its own
files: the \opt{-i} file of \code{wrsamp}, \code{edf2mit}, \code{wav2mit}, \code{a2m} and
\code{ad2m}, the input files of \code{ahaecg2mit}, the \opt{-S} script of \code{xform} and
the signal file that \code{wfdbdesc} measures are read under the input directory, and the
WAV file of \code{mit2wav}, the EDF file of \code{mit2edf}, the master header of
\cmd{snip -m} and the \file{.des}, \file{.ecg}, \file{.key} and \file{.txt} files of
\code{parsescp} are created under the output directory (with
\code{wfdbx\_input\_path} and \code{wfdbx\_output\_path}, \cref{sec:kotlin-library}). Such
a file cannot be read from a URL. In legacy mode, they are read as named, relative to the
current directory, and the output directory, if set, is a prefix of their names.

\paragraph{The launcher.} \file{bin/wfdbx} runs \cmd{java -cp build/libs/wfdbx-all.jar
wfdbx.cli.MainKt TOOL ...}, or the jar named by the environment variable
\code{WFDBX\_KOTLIN\_JAR}. It starts the JVM with options for a short start-up (its first
compiler only, the serial collector, and the class data sharing archive
\file{wfdbx-all.jsa} that \cmd{./gradlew fatJar} makes next to the jar), and moves
\code{JAVA\_TOOL\_OPTIONS} to the command line. A link to \file{bin/wfdbx} named after a
tool runs that tool. Without the launcher, the jar runs the tools directly:

\begin{lstlisting}[style=shell]
java -cp build/libs/wfdbx-all.jar wfdbx.cli.MainKt rdsamp -r 100s -t 10
\end{lstlisting}

\paragraph{The tools as functions.} Each tool is a function of the package
\code{wfdbx.apps}, such as
\begin{lstlisting}[language=Kotlin]
fun Wfdb.rdsampMain(argv: Array<ByteArray>): Int
\end{lstlisting}
where \code{argv} is C's \code{argv}, the program name first, and the result is the exit
status. A tool may end the process (with
\code{exit}, as C programs do), so an app uses the high-level functions or the library
functions instead.

\paragraph{Environment variables.} The tools, like the rest of the library, use the
environment variables of \cref{tab:kotlin-environment}. On the JVM, which cannot change
its own environment, \code{setenv} (in \code{wfdbx.cstd}) keeps values that the library
and the processes it starts see.

\begin{table}[htbp]
\centering
\caption{The environment variables used by the Kotlin library}
\label{tab:kotlin-environment}
\begin{tabularx}{\linewidth}{@{}lL@{}}
\toprule
\tabhead Variable & \tabhead Meaning \\
\midrule
\code{WFDBX\_INPUT\_DIR} & the input directory (a directory, a \code{\textasciitilde}
  path or a URL), when \code{wfdbxSetInputDir} does not set one \\
\code{WFDBX\_OUTPUT\_DIR} & the output directory, when \code{wfdbxSetOutputDir} does not
  set one \\
\code{WFDBX\_NETWORK} & network access is on if it is a nonzero number, when
  \code{wfdbxSetNetwork} does not set it \\
\code{WFDBANNSORT} & \code{0} to keep annotation files in the order written (no sorting
  when they are closed) \\
\code{WFDBGVMODE} & the default getvec mode (\code{setgvmode} with a negative mode) \\
\code{WFDBCAL} & the calibration file read by \code{calopen(null)} (default
  \file{wfdbcal}) \\
\code{WFDB} & the WFDB path, used in legacy mode only \\
\code{WFDB\_FLAC\_CHECK\_MD5} & if set, the MD5 signature of FLAC signal files is
  verified when they are read \\
\code{WFDBX\_KOTLIN\_JAR} & the jar that \file{bin/wfdbx} runs (default
  \file{build/libs/wfdbx-all.jar}) \\
\bottomrule
\end{tabularx}
\end{table}

% ==========================================================================================
\section{Testing}
\label{sec:kotlin-testing}

The unit tests of the Kotlin library are in \file{kotlin/src/commonTest} (run on every
target) and \file{kotlin/src/jvmTest} (JVM only). In the directory \file{kotlin}:

\begin{lstlisting}[style=shell]
./gradlew jvmTest                           # the unit tests, on the JVM
./gradlew compileCommonMainKotlinMetadata   # the common code uses only common APIs
./gradlew jvmTest -Pwfdbx.linux=true linuxX64Test   # also on Linux/Native, with its toolchain
\end{lstlisting}

The unit tests check the library functions with expected values printed by the C
library, the high-level functions, the bulk paths against the per-frame paths, the header
metadata functions against the cases shared by the seven languages (in
\file{tests/metadata}: results, warnings and the headers, byte for byte), the examples of
\file{kotlin/README.md}, and \code{checkRecord} against the reports of the Python
implementation on a corpus of damaged records (this test is skipped, with a message, when
\code{python3} with NumPy is not available). Network access is tested against a local
HTTP server (127.0.0.1). The tests copy \file{tests/data} to a scratch directory, and
never change the current directory.

The cross-language conformance suite compares the output of the tools with that of the
original C programs, case by case: 1135 cases, and 883 in legacy mode, which runs each
tool with \opt{--legacy}. It runs the launcher \file{kotlin/bin/wfdbx}, so the jar is
built first. From the root of the repository:

\begin{lstlisting}[style=shell]
(cd kotlin && ./gradlew fatJar)
PYTHONPATH=$PWD/python/src python3 tests/conformance/harness.py check --impl kotlin -j8
PYTHONPATH=$PWD/python/src python3 tests/conformance/harness.py check --impl kotlin -j8 --legacy
\end{lstlisting}
